\documentclass[11pt,oneside]{article}

\usepackage[T1]{fontenc}
\usepackage[utf8]{inputenc}
\usepackage{lmodern}
\usepackage{microtype}
\usepackage[a4paper,margin=1in,headheight=14pt]{geometry}
\usepackage{amsmath,amssymb,amsthm,mathtools}
\usepackage{booktabs,longtable,array,tabularx,multirow}
\usepackage{enumitem}
\usepackage{xcolor}
\usepackage{fancyhdr}
\usepackage{titlesec}
\usepackage{tikz}
\usetikzlibrary{arrows.meta,positioning,shapes.geometric}
\usepackage{listings}
\usepackage{url}
\usepackage{hyperref}
\usepackage[nameinlink,capitalise,noabbrev]{cleveref}

\definecolor{deepblue}{RGB}{33,72,116}
\definecolor{deepgreen}{RGB}{35,104,75}
\definecolor{darkred}{RGB}{139,44,44}
\definecolor{softgray}{RGB}{244,246,248}
\definecolor{midgray}{RGB}{95,104,116}

\hypersetup{
  colorlinks=true,
  linkcolor=deepblue,
  citecolor=deepgreen,
  urlcolor=deepblue,
  pdftitle={Valuation-Slope Energy and Mobius Discriminants in the Alaoglu--Erdos Problem},
  pdfauthor={Research continuation prepared for Vladimir Reshetnikov}
}

\pagestyle{fancy}
\fancyhf{}
\fancyhead[L]{Alaoglu--Erdős progress report}
\fancyhead[R]{22 August 2026}
\fancyfoot[C]{\thepage}

\titleformat{\section}{\Large\bfseries\color{deepblue}}{\thesection}{0.7em}{}
\titleformat{\subsection}{\large\bfseries\color{deepblue}}{\thesubsection}{0.7em}{}
\titleformat{\subsubsection}{\normalsize\bfseries\color{deepgreen}}{\thesubsubsection}{0.7em}{}

\newtheoremstyle{research}
  {7pt}{7pt}{\itshape}{} {\bfseries\color{deepblue}}{.}{0.5em}{}
\theoremstyle{research}
\newtheorem{theorem}{Theorem}[section]
\newtheorem{proposition}[theorem]{Proposition}
\newtheorem{lemma}[theorem]{Lemma}
\newtheorem{corollary}[theorem]{Corollary}
\newtheorem{conjecture}[theorem]{Target statement}

\theoremstyle{definition}
\newtheorem{definition}[theorem]{Definition}
\newtheorem{remark}[theorem]{Remark}
\newtheorem{example}[theorem]{Example}

\newcommand{\R}{\mathbb R}
\newcommand{\Q}{\mathbb Q}
\newcommand{\Z}{\mathbb Z}
\newcommand{\N}{\mathbb N}
\newcommand{\PP}{\mathcal P}
\newcommand{\E}{\mathcal E}
\newcommand{\TV}{\operatorname{TV}}
\newcommand{\code}[1]{\texttt{\detokenize{#1}}}
\newcommand{\vp}{v_p}
\newcommand{\ord}{\operatorname{ord}}
\newcommand{\rank}{\operatorname{rank}}
\newcommand{\sgn}{\operatorname{sgn}}
\newcommand{\Span}{\operatorname{span}}
\newcommand{\trop}{\mathrm{trop}}
\newcommand{\act}{\mathrm{act}}
\newcommand{\ext}{\mathrm{ext}}
\newcommand{\struct}{\mathrm{str}}
\newcommand{\defeq}{\mathrel{:=}}
\newcommand{\dd}{\,\mathrm d}
\newcommand{\eps}{\varepsilon}
\newcommand{\statusproved}{\textcolor{deepgreen}{\textbf{proved here}}}
\newcommand{\statusinherited}{\textcolor{deepblue}{\textbf{inherited input}}}
\newcommand{\statusopen}{\textcolor{darkred}{\textbf{open}}}

\newenvironment{resultbox}
  {\begin{center}\begin{minipage}{0.94\textwidth}\setlength{\parindent}{0pt}%
   \hrule\vspace{0.45em}}
  {\vspace{0.45em}\hrule\end{minipage}\end{center}}

\lstset{
  basicstyle=\ttfamily\footnotesize,
  columns=fullflexible,
  frame=single,
  backgroundcolor=\color{softgray},
  rulecolor=\color{midgray},
  breaklines=true,
  keepspaces=true,
  literate=
    {·}{{$\cdot$}}1
    {δ}{{$\delta$}}1
    {θ}{{$\theta$}}1
    {ι}{{$\iota$}}1
    {→}{{$\to$}}1
    {↔}{{$\leftrightarrow$}}1
    {∀}{{$\forall$}}1
    {∃}{{$\exists$}}1
    {∑}{{$\sum$}}1
    {∧}{{$\land$}}1
    {∨}{{$\lor$}}1
    {∥}{{$\Vert$}}1
    {⟂}{{$\perp$}}1
    {≠}{{$\ne$}}1
    {≥}{{$\ge$}}1
    {≤}{{$\le$}}1
    {⟨}{{$\langle$}}1
    {⟩}{{$\rangle$}}1
}

\begin{document}

\begin{titlepage}
\centering
\vspace*{1.2cm}
{\Huge\bfseries\color{deepblue}
Valuation-Slope Energy and M\"obius Discriminants\par}
\vspace{0.35cm}
{\LARGE\bfseries in the Alaoglu--Erdős Problem\par}
\vspace{0.8cm}
{\Large A nonduplicative continuation of the unified working report\par}
\vspace{1.3cm}
{\large Research progress report\par}
\vspace{0.2cm}
{\large 22 August 2026\par}
\vfill
\begin{minipage}{0.84\textwidth}
\small
\textbf{Problem.}  Prove that for real $x$,
\[
  2^x,3^x\in\Z \quad\Longrightarrow\quad x\in\Z.
\]

\textbf{Status.}  This continuation does not claim a complete proof.  It proves two
new exact reductions that sharpen the two surviving arithmetic branches of the
attached unified report.  The rank-four branch acquires a quantitative
valuation-slope energy and an explicit integer-wedge gap.  The rank-three branch
acquires a canonical integer M\"obius discriminant, identified exactly with the
determinant of the formal logarithmic quadratic.  In the no-cross-prime
subbranch this leaves only two one-sided families and produces a second,
localized four-exponentials square.
\end{minipage}
\vfill
{\small Prepared as a standalone continuation; inherited results are cited but not reproved.}
\end{titlepage}

\begin{abstract}
Let $\beta>0$ be a hypothetical nonintegral real number for which
$M=2^\beta$ and $N=3^\beta$ are integers, and let
$\theta=\log_2 3$.  The supplied unified report reduces the problem to a
small number of rigid branches but leaves two persistent obstructions: a
rank-four valuation branch and a rank-three M\"obius branch.

This report makes both obstructions discrete.  In the cross-coprime rank-four
branch, the normalized external prime-log factorizations of $M$ and $N$ are
probability distributions $u=(u_p)$ and $v=(v_p)$.  The Gini norm governing the
mixed determinant has the exact first-quadrant formula
\[
 \mathfrak G(u,v)=\frac{2}{15}\left(u+v+
 \frac{(u-v)^2}{\max(u,v)}\right),
\]
with the final quotient defined as zero at $(0,0)$.  Consequently the old
$7/10$ tropical ceiling is exactly
\[
 \frac{\log Q_T^{\trop}}{H_T}\longrightarrow
 \frac{7}{10}-\frac{\E}{20},
 \qquad
 \E=\sum_{p\ge5}\frac{(u_p-v_p)^2}{\max(u_p,v_p)}.
\]
Every nonzero local summand is strict.  More importantly, prime-slope
straddling supplies two primes $p<q$ with a nonzero integer minor
$\omega_{pq}=v_p(M)v_q(N)-v_q(M)v_p(N)$.  Writing
$\kappa=(\log5\log7)/(\log2\log3)$, one obtains
\[
 \TV(u,v)\ge \frac{\kappa}{\beta^2},\qquad
 \E\ge \frac{4\kappa^2}{\beta^4+\kappa\beta^2},
\]
and hence a completely explicit strict deficit below $7/10$.

In the rank-three branch the external valuation vectors lie on one primitive
ray, so
\[
 M=2^{\alpha_2}3^{\alpha_3}W^r,
 \qquad
 N=2^{\gamma_2}3^{\gamma_3}W^s
\]
for an integer $W>1$ coprime to $6$.  Setting
$U=s\alpha_2-r\gamma_2$ and $V=s\alpha_3-r\gamma_3$ gives
\[
 \beta=\frac{U+V\theta}{s-r\theta}.
\]
The integer
\[
 \Delta=-rU-sV
 =r^2\gamma_2+rs(\gamma_3-\alpha_2)-s^2\alpha_3
\]
is nonzero for every nonintegral solution.  It is simultaneously the negative
M\"obius determinant and, up to the factor $-1/4$, the determinant of the
formal quadratic relation among $\log2,\log3,\log W$.  Under the no-cross-prime
conditions $3\nmid M$ and $2\nmid N$, leastness leaves exactly two one-sided
families; both reduce to
\[
 W^{\lvert r\theta-s\rvert}=3^d,
 \qquad d\in\N,
\]
and yield an explicit $2\times2$ four-exponentials square with bases $2,W$.
The report closes with exact completion criteria, failed closure attempts,
and a Lean-oriented formalization plan.
\end{abstract}

\tableofcontents
\newpage

\section{Executive result ledger}

The purpose of this section is to state precisely what has changed since the
unified report and what has not.  Every assertion is assigned one of three
statuses: \statusproved, \statusinherited, or \statusopen.

\begin{longtable}{p{0.25\textwidth}p{0.47\textwidth}p{0.18\textwidth}}
\toprule
Item & Content & Status\\
\midrule
\endfirsthead
\toprule
Item & Content & Status\\
\midrule
\endhead
Least exponent and primitive output & A counterexample yields a least
nonintegral $\beta>0$ with integer outputs $M=2^\beta$, $N=3^\beta$, together
with the semigroup and primitive-depth reductions developed in the supplied
report. & \statusinherited\\
\addlinespace
Mixed determinant asymptotic & For the determinant family of the unified
report, the leading tropical divisor is controlled by a sum of Gini norms of
normalized prime vectors. & \statusinherited\\
\addlinespace
First-quadrant Gini identity & The Gini norm equals a linear term plus the
explicit diagonal-defect term $(u-v)^2/\max(u,v)$. & \statusproved\\
\addlinespace
Exact cross-coprime ceiling & In the cross-coprime branch the tropical leading
constant is exactly $7/10-\E/20$. & \statusproved\\
\addlinespace
Prime-slope straddling & The local errors $v_p(M)\theta-v_p(N)$ have both
signs; therefore two external valuation columns have a positive integer
minor. & \statusproved\\
\addlinespace
Quantitative rank-four gap & The integer minor forces an explicit
$\beta^{-2}$ total-variation gap and a $\beta^{-4}$ energy gap. & \statusproved\\
\addlinespace
Rank-three ray factorization & Collinear external valuations admit a unique
primitive external core $W$ and exponents $r,s$. & \statusinherited\newline\textbf{[elementary]}\\
\addlinespace
M\"obius discriminant & The integer $\Delta=-rU-sV$ is nonzero and is exactly
the determinant controlling both the M\"obius map and the formal quadratic. & \statusproved\\
\addlinespace
No-cross rank-three classification & If $3\nmid M$ and $2\nmid N$, then exactly
one of the two base depths survives, giving two one-sided localized radical
families. & \statusproved\\
\addlinespace
Complete Alaoglu--Erdős proof & Elimination of all rank-four and rank-three
branches. & \statusopen\\
\bottomrule
\end{longtable}

\begin{resultbox}
\textbf{Main advance in one sentence.}
The two continuous-looking obstructions in the unified report now carry
nonzero integer certificates: an exterior valuation minor $\omega_{pq}$ in
rank four, or a M\"obius discriminant $\Delta$ in rank three.
\end{resultbox}

\section{Scope, nonduplication, and current reconnaissance}

\subsection{What is deliberately not repeated}

The supplied manuscript is already a large synthesis of many continuations.
It develops, among other things, minimal-counterexample semigroups,
generalized Vandermonde determinants, min-plus ceilings, Hermite--Chebyshev
interpolation, $p$-adic layers, formal logarithmic quadratics, matrix
coefficient bounds, canonical products, P\'olya--Carlson and Mahler routes,
operator and moment formulations, and an extensive computational frontier.
Repeating those constructions would make this report longer without making it
stronger.  They are therefore compressed below into a short list of inherited
inputs.

The genuinely new material begins with \cref{sec:gini-energy}.  In particular,
the following formulas and deductions were not found in the supplied source:
\begin{enumerate}[label=(N\arabic*),leftmargin=2.7em]
\item the diagonal-defect formula for the first-quadrant Gini norm;
\item the exact identity $7/10-\E/20$ for the cross-coprime tropical constant;
\item the probability-distribution and total-variation interpretation;
\item the signed prime-slope straddling lemma and its integer minor;
\item the explicit $\beta^{-4}$ gap obtained from that minor;
\item the valuation M\"obius determinant $\Delta$ and its equality with the
      formal quadratic determinant;
\item the complete no-cross-prime rank-three zero-pattern classification;
\item the localized four-exponentials square with bases $2$ and $W$.
\end{enumerate}

\subsection{Public claims surrounding the competing system}

The request motivating this continuation mentioned recent public claims about
an experimental system labelled ``Fable 5.''  I checked the publicly available
artifacts on 22 August 2026.  The following is the evidence classification used
here.

\begin{itemize}[leftmargin=2em]
\item Anthropic's official site publicly describes Claude Opus 4.6; the
      ``Fable 5'' label appears on the independent challenge site rather than
      as a separate official model card \cite{AnthropicOpus46,ByClaudeHome}.
\item A public challenge page reports an elliptic curve of rank at least $30$
      and links a paper/source package attributed to Levent Alp\"oge
      \cite{ByClaudeRank30}.  This is a concrete public artifact, although this
      report does not independently re-run the complete arithmetic
      certification.
\item A public Jacobian-conjecture verifier and associated challenge pages are
      available \cite{JacobianVerifier,ByClaudeJacobian}.  Again, their
      mathematical audit is a separate project from the present one.
\item The Hadamard challenge page reports certificates for the twelve formerly
      unresolved orders below $2000$ and links the claimed constructions
      \cite{ByClaudeHadamard}.  I verified the existence of the public index,
      not every matrix certificate.
\end{itemize}

No public preprint, Lean development, or independently inspectable certificate
for a new Alaoglu--Erdős breakthrough was located in this search.  That
negative observation is necessarily scoped to the public sources available at
the search date; the private message described in the request is not directly
verifiable.  None of the mathematics below depends on any competitive claim.

\subsection{Why this continuation attacks invariants rather than methods}

The unified report repeatedly reaches one of two walls:
\begin{enumerate}[label=(\roman*),leftmargin=2em]
\item in rank four, termwise determinant divisibility saturates at a leading
      constant below the archimedean bound;
\item in rank three, the surviving logarithmic relation is a full-rank
      quadratic, equivalently a rational M\"obius relation between the
      hypothetical exponent and $\theta=\log_2 3$.
\end{enumerate}
A new interpolation kernel, a different basis, or another norm estimate tends
to reproduce one of those walls.  The present strategy is therefore to expose
the discrete arithmetic hidden inside each wall.  A nonzero integer is harder
to deform away than an inequality with an unquantified strictness clause.

\section{Inherited setup and notation}
\label{sec:setup}

Assume, for the purpose of reduction, that the conjecture is false.  The
minimal-counterexample machinery in the unified report permits us to choose a
least nonintegral exponent
\[
  \beta>0,
  \qquad M=2^\beta\in\N,
  \qquad N=3^\beta\in\N.
\]
Set
\[
  \theta\defeq\log_2 3.
\]
Unique factorization shows that $\theta\notin\Q$.  Indeed, if
$\theta=a/b$ in lowest terms, then $2^a=3^b$, which is impossible.

For each prime $p$, put
\[
  a_p\defeq v_p(M),\qquad b_p\defeq v_p(N).
\]
We reserve the structural depths
\[
 \alpha_2=v_2(M),\quad \alpha_3=v_3(M),\quad
 \gamma_2=v_2(N),\quad \gamma_3=v_3(N).
\]
The external valuation vectors are
\[
 \mathbf a_{\ext}=(a_p)_{p\ge5},\qquad
 \mathbf b_{\ext}=(b_p)_{p\ge5}.
\]

We use the following inherited facts from the unified corpus.

\begin{description}[style=nextline,leftmargin=2.3em]
\item[H1: leastness and primitive depth.]
A least counterexample satisfies a primitive output condition.  In particular,
the two base depths obey
\[
 \alpha_2=0\quad\text{or}\quad \gamma_3=0.
\]
The stronger common-perfect-power exclusions are not needed below.

\item[H2: rational exponents are already settled.]
If $\beta\in\Q$ and $2^\beta,3^\beta\in\Z$, then $\beta\in\Z$.  This is also
immediate: writing $\beta=c/d$ in lowest terms, integrality of $2^{c/d}$ forces
$d\mid c$ by the $2$-adic valuation.

\item[H3: mixed Gini asymptotic.]
For the mixed determinant family indexed by a scale $T$, the source constructs
a tropical divisor $Q_T^{\trop}$ and an archimedean height $H_T$ such that
\[
 \frac{\log Q_T^{\trop}}{H_T}
 \longrightarrow
 1-\frac38\sum_p\mathfrak G(c_p),
 \tag{H3}
\]
where $c_p$ are normalized prime vectors and $\mathfrak G$ is the simplex
Gini norm defined in \cref{def:gini} below.

\item[H4: rank dichotomy.]
The four logarithmic valuation vectors have only two surviving possibilities.
The external pair is not identically zero.  In rank four, the external valuation
vectors are independent.  In rank three, they are collinear and the
second-iterate kernel is nonzero; equivalently,
$\beta$ lies in a rational M\"obius orbit of $\theta$.

\item[H5: formal quadratic.]
In rank three, unique factorization reduces the logarithmic relation to a
homogeneous quadratic in three independent logarithms.  Baker-type results
exclude rank at most two; the full-rank quadratic is the surviving case.

\item[H6: finite exclusion.]
The accepted finite verification gives $\beta>12$.  This is used only to
simplify the final quantitative energy bound.
\end{description}

The point of the next two parts is to turn H3--H5 into exact arithmetic
statements with nonzero integer witnesses.

\section{The rank-four invariant: valuation-slope energy}
\label{sec:gini-energy}

We first treat the branch in which neither output contains a structural prime:
\begin{equation}
  \gcd(MN,6)=1.
  \label{eq:cross-coprime}
\end{equation}
This is the ``cross-coprime'' branch of the unified report.  Its determinant
analysis produced a universal tropical ceiling of $7/10$, with a qualitative
strictness clause unless all external prime vectors lie on one ray.  The
purpose of this section is to replace that clause by an exact positive energy
and then force that energy away from zero by a nonzero integer minor.

\subsection{Two probability distributions on the external primes}

Under \eqref{eq:cross-coprime}, all prime factors of $M$ and $N$ are at least
$5$.  For every prime $p\ge5$, define
\begin{equation}
 u_p=\frac{a_p\log p}{\beta\log2},
 \qquad
 v_p=\frac{b_p\log p}{\beta\log3}.
 \label{eq:uv-def}
\end{equation}
Only finitely many entries are nonzero, and the two logarithmic
factorizations give
\begin{equation}
 \sum_{p\ge5}u_p=1,
 \qquad
 \sum_{p\ge5}v_p=1.
 \label{eq:uv-probability}
\end{equation}
Thus $u=(u_p)$ and $v=(v_p)$ are probability distributions on the set of
external primes.  This elementary normalization is the key simplification:
the branch defect can be read as a distance between two factorizations rather
than as an unstructured sum over places.

Let
\[
 \Delta_\triangle
 =\{(x_1,x_2)\in\R_{\ge0}^2:x_1+x_2\le1\},
\]
and let $X,X'$ be independent random points uniformly distributed on
$\Delta_\triangle$.

\begin{definition}[Triangular Gini norm]
\label{def:gini}
For $(r,s)\in\R^2$, set
\begin{equation}
 \mathfrak G(r,s)
 =\mathbb E\left|r(X_1-X_1')+s(X_2-X_2')\right|.
 \label{eq:gini-def}
\end{equation}
\end{definition}

The unified report proves that this is a norm and gives a closed formula after
sorting the three vertex values $0,r,s$.  In the first quadrant that formula
has a more revealing form.

\begin{proposition}[Diagonal-defect formula]
\label{prop:gini-first-quadrant}
For $u,v\ge0$, not both zero,
\begin{equation}
 \mathfrak G(u,v)
 =\frac{2}{15}\left(
   u+v+\frac{(u-v)^2}{\max(u,v)}
  \right).
 \label{eq:gini-diagonal-defect}
\end{equation}
At $(u,v)=(0,0)$ the final quotient is interpreted as zero.
\end{proposition}

\begin{proof}
Assume first that $0\le u\le v$.  In the notation of the closed formula in the
unified report, the ordered vertex gaps are
$A_0=u$ and $B_0=v-u$.  Therefore
\[
 \mathfrak G(u,v)
 =\frac{2\bigl(2u^2+3u(v-u)+2(v-u)^2\bigr)}{15v}
 =\frac{2(2v^2-uv+u^2)}{15v}.
\]
Since
\[
 \frac{2v^2-uv+u^2}{v}
 =u+v+\frac{(v-u)^2}{v},
\]
this is \eqref{eq:gini-diagonal-defect}.  The case $v\le u$ follows by
symmetry, and the origin follows by continuity.
\end{proof}

The first term in \eqref{eq:gini-diagonal-defect} is forced by total mass.  All
branch-specific information is carried by the nonnegative second term.

\begin{definition}[Valuation-slope energy]
\label{def:energy}
Define
\begin{equation}
 \E(M,N)
 =\sum_{p\ge5}
   \frac{(u_p-v_p)^2}{\max(u_p,v_p)},
 \label{eq:energy-def}
\end{equation}
again assigning the value zero to a zero coordinate pair.
\end{definition}

\begin{proposition}[Primewise logarithmic form of the energy]
\label{prop:energy-arithmetic-form}
The energy has the exact expression
\begin{equation}
 \E(M,N)
 =\frac1{\beta\log3}
   \sum_{\substack{p\ge5\\a_p+b_p>0}}\log p\,
   \frac{(a_p\theta-b_p)^2}{\max(a_p\theta,b_p)}.
 \label{eq:energy-arithmetic-form}
\end{equation}
\end{proposition}

\begin{proof}
From \eqref{eq:uv-def},
\[
 u_p-v_p
 =\frac{\log p}{\beta\log3}(a_p\theta-b_p),
\]
and
\[
 \max(u_p,v_p)
 =\frac{\log p}{\beta\log3}\max(a_p\theta,b_p).
\]
Substitution gives the result term by term.
\end{proof}

The arithmetic form of the energy shows that the determinant defect is a
weighted sum of local rational-approximation errors to $\theta=\log_2 3$.
It is not an arbitrary analytic correction.

\subsection{The exact cross-coprime determinant constant}

Under \eqref{eq:cross-coprime}, the structural normalized vectors in H3 are
\[
 c_2=(-1,0),\qquad c_3=(0,-1),
\]
while $c_p=(u_p,v_p)$ for $p\ge5$.  Since
$\mathfrak G(1,0)=\mathfrak G(0,1)=4/15$, Proposition
\ref{prop:gini-first-quadrant} and \eqref{eq:uv-probability} give
\begin{align}
 \sum_{p\ge5}\mathfrak G(u_p,v_p)
 &=\frac2{15}\left(
     \sum_{p\ge5}(u_p+v_p)+\E(M,N)
   \right) \\
 &=\frac4{15}+\frac2{15}\E(M,N).
 \label{eq:external-gini-energy}
\end{align}
Hence the full Gini sum in H3 is
\begin{equation}
 \sum_p\mathfrak G(c_p)
 =\frac45+\frac2{15}\E(M,N).
 \label{eq:full-gini-energy}
\end{equation}

\begin{theorem}[Exact energy correction to the $7/10$ ceiling]
\label{thm:exact-cross-coprime-ceiling}
In the cross-coprime branch,
\begin{equation}
 \frac{\log Q_T^{\trop}}{H_T}
 \longrightarrow
 \frac7{10}-\frac1{20}\E(M,N).
 \label{eq:exact-cross-coprime-ceiling}
\end{equation}
Moreover
\begin{equation}
 0<\E(M,N)\le2,
 \qquad
 \frac35\le
 \lim_{T\to\infty}\frac{\log Q_T^{\trop}}{H_T}
 <\frac7{10}.
 \label{eq:energy-range}
\end{equation}
\end{theorem}

\begin{proof}
Substitute \eqref{eq:full-gini-energy} into H3:
\[
 1-\frac38\left(\frac45+\frac2{15}\E\right)
 =\frac7{10}-\frac{\E}{20}.
\]
For the upper bound on the energy, put $d_p=u_p-v_p$ and
$m_p=\max(u_p,v_p)$.  Since $|d_p|\le m_p$,
\[
 \E=\sum_p\frac{d_p^2}{m_p}
 \le\sum_p|d_p|\le2,
\]
the last inequality being the standard $\ell^1$ bound for two probability
distributions.

It remains to prove strict positivity.  If a nonzero coordinate pair had
$u_p=v_p$, then
\[
 a_p\log3=b_p\log2,
\]
so $3^{a_p}=2^{b_p}$.  Unique factorization forces
$a_p=b_p=0$, contrary to nonzero support.  Since $M,N>1$, at least one
external coordinate pair is nonzero; therefore every nonzero summand in
\eqref{eq:energy-def} is positive and $\E>0$.
\end{proof}

The proof identifies the equality case that was left open in the qualitative
triangle-inequality argument of the unified report.  It cannot occur for an
arithmetic output pair: equality would require every external prime to carry
exactly the same normalized mass in $M$ and $N$, which would require a
nontrivial equality between a power of $2$ and a power of $3$ at each prime.

\subsection{Prime-slope straddling and an integer exterior minor}

The positivity of $\E$ is qualitative.  A stronger fact comes from the signed
local slopes
\begin{equation}
 \delta_p=a_p\theta-b_p.
 \label{eq:delta-p}
\end{equation}
Their logarithmically weighted sum vanishes:
\begin{align}
 \sum_{p\ge5}(\log p)\delta_p
 &=\theta\sum_{p\ge5}a_p\log p
   -\sum_{p\ge5}b_p\log p \\
 &=\theta\beta\log2-\beta\log3=0.
 \label{eq:slope-conservation}
\end{align}
No nonzero $\delta_p$ can vanish, by the same unique-factorization argument as
above.

\begin{lemma}[Prime-slope straddling]
\label{lem:prime-slope-straddling}
There are distinct primes $p,q\ge5$ such that
\begin{equation}
 \delta_p>0,
 \qquad
 \delta_q<0.
 \label{eq:straddling-signs}
\end{equation}
For any such ordered pair, the valuation minor
\begin{equation}
 \omega_{pq}
 =a_pb_q-a_qb_p
 \label{eq:omega-def}
\end{equation}
is a positive integer.  In particular, $\omega_{pq}\ge1$.
\end{lemma}

\begin{proof}
If all nonzero $\delta_p$ had the same sign, the finite weighted sum
\eqref{eq:slope-conservation}, whose weights $\log p$ are positive, could not
vanish.  Thus both signs occur.

For a pair satisfying \eqref{eq:straddling-signs},
\begin{align}
 \omega_{pq}
 &=a_p(b_q-a_q\theta)+a_q(a_p\theta-b_p).
 \label{eq:omega-positive-decomposition}
\end{align}
The first term is strictly positive because $a_p>0$ and
$b_q-a_q\theta>0$; the second is nonnegative.  Hence
$\omega_{pq}>0$.  It is an integer by definition.
\end{proof}

This minor is the promised discrete witness for rank four.  It proves directly
that the external valuation columns are not proportional.  It also measures a
normalized area.

\begin{proposition}[Normalized wedge identity]
\label{prop:normalized-wedge}
For any primes $p,q\ge5$,
\begin{equation}
 w_{pq}\defeq u_pv_q-u_qv_p
 =\frac{\log p\log q}{\beta^2\log2\log3}\,\omega_{pq}.
 \label{eq:normalized-wedge}
\end{equation}
For the straddling pair of Lemma~\ref{lem:prime-slope-straddling},
$w_{pq}>0$ and
\begin{equation}
 w_{pq}\ge
 \frac{\log5\log7}{\beta^2\log2\log3}.
 \label{eq:wedge-lower}
\end{equation}
\end{proposition}

\begin{proof}
Expand \eqref{eq:uv-def}; the common factors give
\eqref{eq:normalized-wedge}.  For a straddling pair,
$\omega_{pq}\ge1$.  The primes are distinct and at least $5$, so
$(\log p)(\log q)\ge(\log5)(\log7)$.
\end{proof}

\subsection{Total variation and the explicit $\beta^{-4}$ gap}

Let
\begin{equation}
 \TV(u,v)=\frac12\sum_{p\ge5}|u_p-v_p|.
 \label{eq:tv-def}
\end{equation}
Because $u$ and $v$ have the same total mass,
\begin{equation}
 \TV(u,v)
 =\sum_{u_p>v_p}(u_p-v_p)
 =\sum_{u_p<v_p}(v_p-u_p).
 \label{eq:tv-positive-negative}
\end{equation}

\begin{lemma}[Energy--variation comparison]
\label{lem:energy-tv}
For any two finitely supported probability distributions $u,v$,
\begin{equation}
 \frac{4\TV(u,v)^2}{1+\TV(u,v)}
 \le \sum_p\frac{(u_p-v_p)^2}{\max(u_p,v_p)}
 \le 2\TV(u,v).
 \label{eq:energy-tv}
\end{equation}
\end{lemma}

\begin{proof}
Write $d_p=u_p-v_p$ and $m_p=\max(u_p,v_p)$.  Then
\[
 \sum_p m_p
 =\frac12\sum_p(u_p+v_p+|d_p|)
 =1+\TV(u,v).
\]
Cauchy--Schwarz gives
\[
 (2\TV)^2
 =\left(\sum_p|d_p|\right)^2
 \le\left(\sum_p\frac{d_p^2}{m_p}\right)
    \left(\sum_pm_p\right),
\]
which is the lower bound.  Since $|d_p|\le m_p$, one has
$d_p^2/m_p\le|d_p|$ termwise, giving the upper bound.
\end{proof}

\begin{lemma}[A straddling wedge is bounded by total variation]
\label{lem:wedge-tv}
If $u_p>v_p$ and $u_q<v_q$, then
\begin{equation}
 0<u_pv_q-u_qv_p\le\TV(u,v).
 \label{eq:wedge-tv}
\end{equation}
\end{lemma}

\begin{proof}
Put $d_p=u_p-v_p>0$ and $e_q=v_q-u_q>0$.  Then
\[
 u_pv_q-u_qv_p=v_qd_p+v_pe_q.
\]
By \eqref{eq:tv-positive-negative}, $d_p,e_q\le\TV(u,v)$, while
$v_p+v_q\le1$.  Hence the displayed quantity is at most
$\TV(u,v)(v_p+v_q)\le\TV(u,v)$.
\end{proof}

Define the absolute constant
\begin{equation}
 \kappa
 =\frac{\log5\log7}{\log2\log3}
 =4.1127006240\ldots .
 \label{eq:kappa}
\end{equation}

\begin{theorem}[Explicit arithmetic separation]
\label{thm:explicit-energy-gap}
Every hypothetical cross-coprime counterexample satisfies
\begin{equation}
 \TV(u,v)\ge\frac{\kappa}{\beta^2},
 \qquad
 \|u-v\|_1\ge\frac{2\kappa}{\beta^2},
 \label{eq:tv-lower}
\end{equation}
and
\begin{equation}
 \E(M,N)
 \ge
 \frac{4\kappa^2}{\beta^4+\kappa\beta^2}.
 \label{eq:energy-lower}
\end{equation}
Consequently
\begin{equation}
 \lim_{T\to\infty}\frac{\log Q_T^{\trop}}{H_T}
 \le
 \frac7{10}
 -\frac{\kappa^2}{5(\beta^4+\kappa\beta^2)}.
 \label{eq:quantitative-ceiling}
\end{equation}
Since the inherited finite exclusion gives $\beta>12$, one may use the simpler
bound
\begin{equation}
 \lim_{T\to\infty}\frac{\log Q_T^{\trop}}{H_T}
 \le
 \frac7{10}-\frac{\kappa^2}{10\beta^4}
 =\frac7{10}-\frac{1.6914306423\ldots}{\beta^4}.
 \label{eq:simple-quantitative-ceiling}
\end{equation}
\end{theorem}

\begin{proof}
Choose a straddling pair from Lemma~\ref{lem:prime-slope-straddling}.
Proposition~\ref{prop:normalized-wedge} and Lemma~\ref{lem:wedge-tv} give
\[
 \TV(u,v)\ge w_{pq}\ge\frac\kappa{\beta^2},
\]
which proves \eqref{eq:tv-lower}.  The function
$t\mapsto4t^2/(1+t)$ is increasing for $t\ge0$, so
Lemma~\ref{lem:energy-tv} gives
\[
 \E\ge
 \frac{4(\kappa/\beta^2)^2}{1+\kappa/\beta^2}
 =\frac{4\kappa^2}{\beta^4+\kappa\beta^2}.
\]
Insert this in \eqref{eq:exact-cross-coprime-ceiling} to obtain
\eqref{eq:quantitative-ceiling}.

Finally, $\kappa/\beta^2<1$ when $\beta>12$, and therefore
$1+\kappa/\beta^2<2$.  Thus
$\E\ge2\kappa^2/\beta^4$, which gives
\eqref{eq:simple-quantitative-ceiling}.
\end{proof}

\begin{corollary}[Cross-coprime rank is necessarily four]
\label{cor:cross-coprime-rank-four}
A hypothetical cross-coprime output pair cannot lie in the rank-three
external-ray branch.  Its external valuation matrix has a nonzero
$2\times2$ minor, and its union of external prime supports contains at least
two primes.
\end{corollary}

\begin{proof}
Lemma~\ref{lem:prime-slope-straddling} produces distinct primes $p,q$ with
$\omega_{pq}\ne0$.  Hence the two external valuation vectors are linearly
independent.
\end{proof}

\begin{remark}[What the quantitative deficit means]
The new deficit points in the diagnostically important but unfavorable
direction: it makes the guaranteed termwise divisor smaller than $7/10$, not
larger.  Therefore it does not by itself prove the conjecture.  It does,
however, remove the last possible equality case from the min-plus analysis and
measures exactly how much additional $p$-adic cancellation a successful
rank-four determinant would have to recover; see
\cref{sec:closure-attempts}.
\end{remark}

\section{The rank-three invariant: a valuation M\"obius discriminant}
\label{sec:mobius}

\subsection{Primitive external-core factorization}

We now enter the other surviving branch H4.  Suppose the external valuation
vectors are collinear over $\Q$.  Since they are finitely supported
nonnegative integer vectors, there is a unique primitive vector
$g=(g_p)_{p\ge5}$ on their common ray and unique $r,s\in\N_0$ such that
\[
 a_p=rg_p,
 \qquad
 b_p=sg_p
 \qquad(p\ge5).
\]
Here ``primitive'' means that the greatest common divisor of the nonzero
coordinates of $g$ is one.  Define the external core
\begin{equation}
 W=\prod_{p\ge5}p^{g_p}.
 \label{eq:W-core}
\end{equation}
Then $W>1$, $\gcd(W,6)=1$, and
\begin{equation}
 \boxed{
 M=2^{\alpha_2}3^{\alpha_3}W^r,
 \qquad
 N=2^{\gamma_2}3^{\gamma_3}W^s.}
 \label{eq:rank3-factorization}
\end{equation}
At least one of $r,s$ is nonzero.  The cases $r=0$ or $s=0$ are allowed and
will be classified explicitly below.

The output identities $M=2^\beta$ and $N=3^\beta$ imply
\begin{align}
 \beta\log2
 &=\alpha_2\log2+\alpha_3\log3+r\log W,
 \label{eq:M-log}\\
 \beta\log3
 &=\gamma_2\log2+\gamma_3\log3+s\log W.
 \label{eq:N-log}
\end{align}
Multiply \eqref{eq:M-log} by $s$, multiply \eqref{eq:N-log} by $r$, and
eliminate $\log W$.  After dividing by $\log2$, one obtains
\begin{equation}
 \beta(s-r\theta)=U+V\theta,
 \label{eq:mobius-linear}
\end{equation}
where
\begin{equation}
 U=s\alpha_2-r\gamma_2,
 \qquad
 V=s\alpha_3-r\gamma_3.
 \label{eq:UV}
\end{equation}
Since $\theta$ is irrational and $(r,s)\ne(0,0)$,
\begin{equation}
 s-r\theta\ne0.
 \label{eq:denominator-nonzero}
\end{equation}
Thus
\begin{equation}
 \boxed{
 \beta=f(\theta),
 \qquad
 f(t)=\frac{U+Vt}{s-rt}.}
 \label{eq:mobius-beta}
\end{equation}
This recovers the M\"obius orbit of H4, but now with coefficients written
directly in terms of the valuation depths.

\subsection{The integer discriminant}

\begin{definition}[Valuation M\"obius discriminant]
\label{def:Delta}
For a rank-three factorization \eqref{eq:rank3-factorization}, define
\begin{align}
 \Delta
 &\defeq-rU-sV
 \label{eq:Delta-UV}\\
 &=r^2\gamma_2+rs(\gamma_3-\alpha_2)-s^2\alpha_3.
 \label{eq:Delta-depths}
\end{align}
\end{definition}

The name is justified by the coefficient matrix
\begin{equation}
 B=\begin{pmatrix}
      V & U\\
      -r&s
    \end{pmatrix},
 \qquad
 f(t)=\frac{Vt+U}{-rt+s},
 \qquad
 \det B=sV+rU=-\Delta.
 \label{eq:mobius-matrix}
\end{equation}

\begin{theorem}[Nonvanishing of the M\"obius discriminant]
\label{thm:Delta-nonzero}
For every nonintegral rank-three solution,
\begin{equation}
 \boxed{\Delta\in\Z\setminus\{0\}.}
 \label{eq:Delta-nonzero}
\end{equation}
In particular $|\Delta|\ge1$, and the M\"obius map in
\eqref{eq:mobius-beta} is nonconstant.
\end{theorem}

\begin{proof}
Suppose $\Delta=0$.  Then the two rows $(V,U)$ and $(-r,s)$ of $B$ are
linearly dependent over $\Q$.  Since $(-r,s)\ne(0,0)$, there is
$\lambda\in\Q$ such that
\[
 (V,U)=\lambda(-r,s).
\]
Consequently
\[
 U+V\theta=\lambda(s-r\theta),
\]
and \eqref{eq:mobius-linear} gives $\beta=\lambda\in\Q$.  By H2, $\beta$ is
an integer, contradicting the assumed counterexample.  Thus $\Delta\ne0$.
\end{proof}

\begin{corollary}[M\"obius transversality]
\label{cor:mobius-derivative}
The derivative of the rank-three parametrization is
\begin{equation}
 f'(t)=-\frac{\Delta}{(s-rt)^2},
 \qquad
 |f'(\theta)|=\frac{|\Delta|}{|s-r\theta|^2}.
 \label{eq:mobius-derivative}
\end{equation}
Thus the M\"obius branch is separated from a constant rational map by the
integer gap $|\Delta|\ge1$.
\end{corollary}

\begin{proof}
Differentiate \eqref{eq:mobius-beta}, or use the determinant formula for a
linear fractional transformation.
\end{proof}

\subsection{Identification with the formal logarithmic quadratic}

The formal relation behind the rank-three branch becomes transparent in the
basis
\[
 X=\log2,
 \qquad Y=\log3,
 \qquad Z=\log W.
\]
The equality $(\log M)(\log3)=(\log N)(\log2)$ is
\begin{equation}
 \Phi(X,Y,Z)=0,
 \label{eq:Phi-zero}
\end{equation}
where
\begin{equation}
 \Phi(X,Y,Z)
 =\gamma_2X^2+(\gamma_3-\alpha_2)XY-\alpha_3Y^2
   +sXZ-rYZ.
 \label{eq:formal-q}
\end{equation}
Its symmetric coefficient matrix is
\begin{equation}
 S=\begin{pmatrix}
 \gamma_2 &(\gamma_3-\alpha_2)/2&s/2\\
 (\gamma_3-\alpha_2)/2&-\alpha_3&-r/2\\
 s/2&-r/2&0
 \end{pmatrix}.
 \label{eq:quadratic-matrix}
\end{equation}

\begin{theorem}[The two rank-three determinants are the same integer]
\label{thm:determinant-bridge}
For the matrix in \eqref{eq:quadratic-matrix},
\begin{equation}
 \boxed{\det S=-\frac{\Delta}{4}.}
 \label{eq:detS}
\end{equation}
Hence the following are equivalent:
\begin{enumerate}[label=(\roman*),leftmargin=2em]
\item the M\"obius transformation $f$ is nonconstant;
\item $\Delta\ne0$;
\item the formal quadratic $\Phi$ has rank three.
\end{enumerate}
For a hypothetical nonintegral solution all three conditions hold.
\end{theorem}

\begin{proof}
Expansion along the final row or direct determinant calculation gives
\begin{align*}
 4\det S
 &=-r^2\gamma_2-rs(\gamma_3-\alpha_2)+s^2\alpha_3\\
 &=-\Delta.
\end{align*}
The equivalences follow from \eqref{eq:mobius-matrix} and the criterion
$\det S\ne0$ for a ternary quadratic form to have rank three.  Nonvanishing
for a counterexample is \cref{thm:Delta-nonzero}.
\end{proof}

\begin{resultbox}
\textbf{Conceptual bridge.}
The ``tangency determinant'' of the formal quadratic analysis and the
``PGL$_2(\Q)$ determinant'' of the second-iterate-kernel analysis are not two
parallel obstructions.  After primitive external-core factorization, they are
literally the same integer $\Delta$, up to the harmless factor $-1/4$.
\end{resultbox}

\subsection{Immediate zero-pattern consequences}

The explicit formula \eqref{eq:Delta-depths} resolves several degenerate
support patterns without transcendence theory.

\begin{proposition}[One output has no external core]
\label{prop:one-external-zero}
In the rank-three branch:
\begin{enumerate}[label=(\alph*),leftmargin=2em]
\item If $r=0$ and $s>0$, then $\alpha_3>0$ and
      $\Delta=-s^2\alpha_3<0$.
\item If $s=0$ and $r>0$, then $\gamma_2>0$ and
      $\Delta=r^2\gamma_2>0$.
\end{enumerate}
Equivalently, if $M$ has no external prime factor then it must contain the
cross-prime $3$, and if $N$ has no external prime factor then it must contain
the cross-prime $2$.
\end{proposition}

\begin{proof}
Substitute $r=0$ or $s=0$ into \eqref{eq:Delta-depths} and use
\cref{thm:Delta-nonzero}.
\end{proof}

\begin{corollary}[Cross-coprime rank three is impossible]
\label{cor:cross-coprime-no-rank3}
If $\alpha_2=\alpha_3=\gamma_2=\gamma_3=0$, then $\Delta=0$.  Therefore the
cross-coprime branch cannot have rank three.
\end{corollary}

\begin{proof}
Immediate from \eqref{eq:Delta-depths}.  This is a second proof of
\cref{cor:cross-coprime-rank-four}, now from the rank-three side.
\end{proof}

\begin{proposition}[Height bounds for the discrete certificate]
\label{prop:Delta-height}
Every rank-three factorization satisfies
\begin{align}
 r&\le\frac{\beta\log2}{\log5},&
 s&\le\frac{\beta\log3}{\log5},
 \label{eq:rs-bounds}\\
 0\le\alpha_2&\le\beta,&
 0\le\alpha_3&\le\frac\beta\theta,&
 0\le\gamma_2&\le\beta\theta,&
 0\le\gamma_3&\le\beta.
 \label{eq:depth-bounds}
\end{align}
Consequently
\begin{equation}
 1\le|\Delta|
 \le \frac{3\log2\,\log3}{(\log5)^2}\,\beta^3
 =0.8819474591\ldots\,\beta^3.
 \label{eq:Delta-height-bound}
\end{equation}
\end{proposition}

\begin{proof}
Each nonnegative summand in the logarithm of $M$ or $N$ is bounded by the
whole logarithm.  Since $W\ge5$, equations \eqref{eq:M-log} and
\eqref{eq:N-log} give \eqref{eq:rs-bounds}; the four structural estimates are
similarly immediate.  In \eqref{eq:Delta-depths}, use
$|\gamma_3-\alpha_2|\le\beta$.  Each of the three terms is then bounded by
\[
 \frac{\log2\,\log3}{(\log5)^2}\,\beta^3.
\]
The lower bound is \cref{thm:Delta-nonzero}.
\end{proof}

\begin{remark}
The upper bound is not intended as a contradiction.  It makes the rank-three
search finite at every bounded $\beta$ and shows that the relevant PGL matrix
has polynomial, rather than exponential, coefficient height in $\beta$.  This
is useful for a proof-producing enumeration or a Lean certificate checker.
\end{remark}

\subsection{Complete no-cross-prime classification}
\label{sec:no-cross}

Now impose the no-cross-prime conditions
\begin{equation}
 3\nmid M,
 \qquad
 2\nmid N,
 \qquad\text{equivalently}\qquad
 \alpha_3=\gamma_2=0.
 \label{eq:no-cross}
\end{equation}
These conditions are weaker than full cross-coprimality: $M$ may still contain
$2$, $N$ may still contain $3$, and both may contain the external core $W$.
Under \eqref{eq:no-cross},
\begin{equation}
 \Delta=rs(\gamma_3-\alpha_2).
 \label{eq:Delta-no-cross}
\end{equation}

\begin{theorem}[Two one-sided rank-three families]
\label{thm:no-cross-classification}
Assume rank three, leastness H1, and \eqref{eq:no-cross}.  Then $r,s>0$, and
exactly one of the following two mutually exclusive alternatives holds.

\medskip
\noindent\textbf{Family $+$:}
\begin{equation}
 \alpha_2=0,
 \quad \gamma_3=d>0,
 \quad M=W^r,
 \quad N=3^dW^s,
 \quad r\theta>s,
 \label{eq:family-plus}
\end{equation}
with
\begin{equation}
 \beta=\frac{rd\theta}{r\theta-s},
 \qquad
 \theta-\frac sr=\frac{d\theta}{\beta},
 \qquad
 \Delta=rsd>0.
 \label{eq:family-plus-formulas}
\end{equation}

\medskip
\noindent\textbf{Family $-$:}
\begin{equation}
 \gamma_3=0,
 \quad \alpha_2=d>0,
 \quad M=2^dW^r,
 \quad N=W^s,
 \quad s>r\theta,
 \label{eq:family-minus}
\end{equation}
with
\begin{equation}
 \beta=\frac{sd}{s-r\theta},
 \qquad
 1-\frac{r\theta}{s}=\frac d\beta,
 \qquad
 \Delta=-rsd<0.
 \label{eq:family-minus-formulas}
\end{equation}
\end{theorem}

\begin{proof}
If $r=0$, then \cref{prop:one-external-zero}(a) forces
$\alpha_3>0$, contrary to \eqref{eq:no-cross}; similarly $s=0$ is impossible.
Thus $r,s>0$.

Leastness H1 says $\alpha_2=0$ or $\gamma_3=0$.  Equation
\eqref{eq:Delta-no-cross} and \cref{thm:Delta-nonzero} say
$\gamma_3\ne\alpha_2$.  Therefore exactly one of $\alpha_2,\gamma_3$ is zero
and the other is a positive integer $d$.

If $\alpha_2=0$ and $\gamma_3=d$, then \eqref{eq:UV} gives
$U=0$ and $V=-rd$.  Formula \eqref{eq:mobius-beta} becomes
\[
 \beta=\frac{-rd\theta}{s-r\theta}
 =\frac{rd\theta}{r\theta-s}.
\]
Positivity of $\beta$ forces $r\theta>s$, and rearrangement gives the rational
approximation identity in \eqref{eq:family-plus-formulas}.  The value of
$\Delta$ follows from \eqref{eq:Delta-no-cross}.

The second case is symmetric: $U=sd$, $V=0$, so
$\beta=sd/(s-r\theta)$, positivity forces $s>r\theta$, and rearrangement gives
\eqref{eq:family-minus-formulas}.
\end{proof}

\begin{corollary}[Localized radical equation]
\label{cor:localized-radical}
Both alternatives in \cref{thm:no-cross-classification} imply the single form
\begin{equation}
 \boxed{W^{|r\theta-s|}=3^d,}
 \qquad
 W>1,
 \quad \gcd(W,6)=1,
 \quad r,s,d\in\N.
 \label{eq:localized-radical}
\end{equation}
Conversely, together with the appropriate sign of $r\theta-s$, the equation
reconstructs the corresponding factorization in
\eqref{eq:family-plus} or \eqref{eq:family-minus}.
\end{corollary}

\begin{proof}
In Family $+$, compare $N=3^\beta$ with
$N=3^dW^s$ and use $M=W^r=2^\beta$:
\[
 3^dW^s=3^\beta=(2^\beta)^\theta=(W^r)^\theta,
\]
so $W^{r\theta-s}=3^d$.  In Family $-$, compare
$M=2^\beta=2^dW^r$ with $N=W^s=3^\beta$; taking logarithms and using
$\theta\log2=\log3$ gives $(s-r\theta)\log W=d\log3$.  The converse is the
same algebra in reverse.
\end{proof}

\begin{remark}[A precise rational-approximation reading]
Family $+$ places $s/r$ below $\theta$, with exact relative error
$d/\beta$.  Family $-$ places $r\theta/s$ below $1$, again with exact error
$d/\beta$.  These are not exceptionally strong approximations: when
$r,s\asymp\beta$, the error is naturally of order $1/\beta$.  Therefore a
standard irrationality measure for $\theta$ does not eliminate the families.
The useful new feature is the exact integer parameter $d=|\Delta|/(rs)$ and
the localized radical identity \eqref{eq:localized-radical}.
\end{remark}

\subsection{The localized four-exponentials square}

Equation \eqref{eq:localized-radical} produces a second four-exponentials
configuration, distinct from the original square with bases $2,3$ and
exponents $1,\beta$.

\begin{theorem}[A symmetric localized square]
\label{thm:localized-square}
Assume \eqref{eq:localized-radical} and define
\begin{equation}
 \eta=\frac{\log3}{\log W}
      =\frac{|r\theta-s|}{d}>0.
 \label{eq:eta}
\end{equation}
Then
\begin{enumerate}[label=(\roman*),leftmargin=2em]
\item $\log2$ and $\log W$ are linearly independent over $\Q$;
\item $\theta$ and $\eta$ are linearly independent over $\Q$;
\item all four numbers
\begin{equation}
 2^\theta,
 \quad 2^\eta,
 \quad W^\theta,
 \quad W^\eta
 \label{eq:four-localized}
\end{equation}
are algebraic; in fact $2^\theta=W^\eta=3$ and the other two are explicit
rational radicals.
\end{enumerate}
Thus a no-cross-prime rank-three counterexample gives an explicit
$2\times2$ instance forbidden by the Four Exponentials Conjecture.
\end{theorem}

\begin{proof}
Because $W$ is an integer coprime to $2$ and greater than one, no nonzero
integer powers of $2$ and $W$ coincide; hence their logarithms are
$\Q$-independent.

Let $\sigma=\sgn(r\theta-s)\in\{+1,-1\}$.  Then
$d\eta=\sigma(r\theta-s)$.  If
$q\theta+q'\eta=0$ with $q,q'\in\Q$, substitution gives
\[
 \left(q+\frac{\sigma rq'}d\right)\theta
 =\frac{\sigma sq'}d.
\]
Since $s>0$ and $\theta\notin\Q$, both rational coefficients must vanish;
therefore $q'=q=0$.

The diagonal identities are immediate.  Moreover
\[
 2^\eta
 =\left(\frac{3^r}{2^s}\right)^{\sigma/d},
\]
which is algebraic.  If $\sigma=+1$, then
$W^{r\theta-s}=3^d$ gives
\[
 W^\theta=(3^dW^s)^{1/r};
\]
if $\sigma=-1$, then $W^{s-r\theta}=3^d$ gives
\[
 W^\theta=(W^s/3^d)^{1/r}.
\]
Both are algebraic.
\end{proof}

\begin{remark}[Why six exponentials still does not fire]
The localized square supplies two independent base logarithms and two
independent exponents, exactly the open $2\times2$ threshold.  Every additional
base obtained multiplicatively from $2$ and $W$ has logarithm in the same
two-dimensional $\Q$-space, so it cannot provide the third independent base
needed by the Six Exponentials Theorem.  The natural candidate $3$ has
$3^\theta$ and $3^\eta$ uncontrolled.  The reduction is therefore sharp rather
than cosmetic.
\end{remark}


\section{A unified exterior--M\"obius dichotomy}
\label{sec:branch-dichotomy}

The two preceding sections look rather different.  The rank-four branch is
controlled by a family of local $2\times2$ valuation minors, whereas the
rank-three branch is controlled by one global determinant $\Delta$.  They are
in fact the two complementary outputs of the same exterior-algebra test.

Let
\[
 \mathcal P_{\ext}=\{p:p\ge5\text{ prime}\},
 \qquad
 L_{\ext}=\bigoplus_{p\in\mathcal P_{\ext}}\Z e_p,
\]
and encode the external valuations by
\begin{equation}
 A=\sum_{p\ge5}a_pe_p,
 \qquad
 B=\sum_{p\ge5}b_pe_p.
 \label{eq:external-vectors}
\end{equation}
Their exterior product is
\begin{equation}
 \Omega(M,N)=A\wedge B
 =\sum_{p<q}(a_pb_q-a_qb_p)e_p\wedge e_q
 =\sum_{p<q}\omega_{pq}e_p\wedge e_q.
 \label{eq:omega-global}
\end{equation}
It is an integral, finitely supported bivector.

\begin{theorem}[Exterior--M\"obius dichotomy]
\label{thm:exterior-mobius-dichotomy}
For a hypothetical least nonintegral exponent $\beta$, exactly one of the
following alternatives occurs.

\smallskip
\noindent\textbf{(R4) Exterior branch.}
$\Omega(M,N)\ne0$.  Then the four positive algebraic numbers
$2,3,M,N$ have multiplicative rank four, and at least one valuation minor
$\omega_{pq}$ is a nonzero integer.  In the cross-coprime subbranch,
Lemma~\ref{lem:prime-slope-straddling} produces an orientation with
$\omega_{pq}\ge1$, and Theorem~\ref{thm:explicit-energy-gap} gives the explicit
energy gap.

\smallskip
\noindent\textbf{(R3) M\"obius branch.}
$\Omega(M,N)=0$.  Then the external valuation vectors lie on a unique
primitive ray, so there are a primitive integer $W>1$, coprime to $6$, and
$r,s\in\Z_{\ge0}$, not both zero, such that
\[
 M=2^{\alpha_2}3^{\alpha_3}W^r,
 \qquad
 N=2^{\gamma_2}3^{\gamma_3}W^s.
\]
For these uniquely normalized data, the M\"obius determinant
$\Delta=-rU-sV$ of \eqref{eq:Delta-depths} is a nonzero integer.  The four
numbers $2,3,M,N$ have multiplicative rank three.
\end{theorem}

\begin{proof}
By H4, the external valuation vectors are not both zero and their rational
rank is either one or two.  Their rank is two exactly when
$A\wedge B\ne0$, equivalently when at least one $2\times2$ minor
$\omega_{pq}$ is nonzero.  Since the structural primes $2$ and $3$ already
supply two independent valuation directions, this is precisely the
multiplicative-rank-four branch.

If $A\wedge B=0$, the pair $A,B$ is dependent and not both zero.
Because their coordinates are nonnegative integers, the nonzero member or
members determine a unique
primitive nonnegative integral vector
$C=\sum c_pe_p$ and unique $r,s\in\Z_{\ge0}$, not both zero, such that
$A=rC$ and $B=sC$.  Put $W=\prod p^{c_p}$.  This gives the displayed rank-three
factorization and makes $W$ primitive.  Proposition~\ref{thm:Delta-nonzero}
then gives $\Delta\ne0$.  The alternatives are mutually exclusive because
$\Omega$ cannot be both zero and nonzero.
\end{proof}

\begin{definition}[Discrete branch certificates]
\label{def:branch-certificates}
In branch (R4), define
\begin{equation}
 I_4=\gcd_{p<q}|\omega_{pq}|,
 \label{eq:I4}
\end{equation}
where zero minors do not affect the gcd.  In branch (R3), define
\begin{equation}
 I_3=|\Delta|.
 \label{eq:I3}
\end{equation}
Then the active branch certificate is always a positive integer.
\end{definition}

The certificate $I_4$ is local-to-global: it is assembled from pairs of
external primes.  The certificate $I_3$ is global-to-local: after all external
minors vanish, it measures how the structural $(2,3)$ directions twist the
single external ray.  This is why the two invariants do not compete.  The
second becomes visible only after the first has vanished.

\begin{center}
\fbox{\parbox{0.72\textwidth}{\centering
hypothetical least counterexample with external valuation vectors $A,B$}}

\vspace{0.6em}
$\Downarrow\quad$ test $\Omega=A\wedge B$ $\quad\Downarrow$
\vspace{0.6em}

\noindent
\begin{minipage}[t]{0.45\textwidth}
\centering
\fbox{\parbox{0.90\linewidth}{\centering
$\Omega\ne0$: \textbf{rank four}\\
some $\omega_{pq}\in\Z\setminus\{0\}$\\
cross-coprime: $\E>0$ and an explicit gap}}
\end{minipage}\hfill
\begin{minipage}[t]{0.45\textwidth}
\centering
\fbox{\parbox{0.90\linewidth}{\centering
$\Omega=0$: \textbf{rank three}\\
primitive external core $W$\\
$\Delta=\det B_{\mathrm{Mob}}\in\Z\setminus\{0\}$}}
\end{minipage}
\end{center}

\begin{table}[htbp]
\centering
\caption{Complementary arithmetic data in the two multiplicative-rank
branches.}
\label{tab:branch-comparison}
\small
\begin{tabularx}{\textwidth}{@{}p{0.20\textwidth}XX@{}}
\toprule
 & \textbf{rank four} & \textbf{rank three} \\
\midrule
External geometry
 & two-dimensional; $A\wedge B\ne0$
 & one-dimensional; $A\wedge B=0$ \\
Discrete witness
 & a local minor $\omega_{pq}=a_pb_q-a_qb_p$
 & the global M\"obius determinant $\Delta=-rU-sV$ \\
Integrality gap
 & $|\omega_{pq}|\ge1$
 & $|\Delta|\ge1$ \\
Analytic avatar
 & normalized area $u_pv_q-u_qv_p$ and energy $\E$
 & nondegenerate conic $\Phi(X,Y,Z)=0$ with $\det S=-\Delta/4$ \\
Four-\linebreak exponentials wall
 & the original square $(2,3;\,M,N)$
 & a localized square $(2,W;\,\theta,\eta)$ in the no-cross-prime subbranch \\
Current missing input
 & wedge-sensitive cancellation or a stronger determinant
 & exclusion of the corresponding integral special radical \\
\bottomrule
\end{tabularx}
\end{table}

\subsection{An exact one-radical reduction of the rank-three branch}

The conic identity can be solved for the normalized logarithm of the primitive
core.  This gives a compact arithmetic target that is useful both for proof
search and for formalization.

Let
\begin{equation}
 \rho_*=\frac{
  \gamma_2+(\gamma_3-\alpha_2)\theta-\alpha_3\theta^2
 }{r\theta-s},
 \qquad
 \beta_*=\alpha_2+\alpha_3\theta+r\rho_*.
 \label{eq:rho-beta-star}
\end{equation}
The denominator is nonzero by \eqref{eq:denominator-nonzero}.

\begin{theorem}[One-radical rank-three reduction]
\label{thm:rank3-special-radical}
Fix nonnegative integers
$\alpha_2,\alpha_3,\gamma_2,\gamma_3,r,s$ with $(r,s)\ne(0,0)$,
$r\theta\ne s$, and $\Delta\ne0$.  Assume $\rho_*>0$ and $\beta_*>0$.
Then the following are equivalent.
\begin{enumerate}[label=(\roman*)]
\item There is an integer $W>1$, coprime to $6$, such that
\[
 2^{\beta_*}=2^{\alpha_2}3^{\alpha_3}W^r,
 \qquad
 3^{\beta_*}=2^{\gamma_2}3^{\gamma_3}W^s.
\]
\item The single special radical
\begin{equation}
 W_*=2^{\rho_*}
 =2^{\frac{
  \gamma_2+(\gamma_3-\alpha_2)\theta-\alpha_3\theta^2
 }{r\theta-s}}
 \label{eq:rank3-special-radical}
\end{equation}
is an integer $>1$ coprime to $6$.
\end{enumerate}
When these conditions hold, $W=W_*$ uniquely.  When the data arise from an
existing rank-three solution, its exponent is $\beta_*$.
\end{theorem}

\begin{proof}
If (i) holds, take base-$2$ logarithms and put
$\rho=\log W/\log2$.  Eliminating $\beta$ from
\[
 \beta=\alpha_2+\alpha_3\theta+r\rho,
 \qquad
 \beta\theta=\gamma_2+\gamma_3\theta+s\rho
\]
gives $\rho=\rho_*$, so $W=2^{\rho_*}$.

Conversely, assume (ii) and set $W=W_*$.  The definition of $\beta_*$ gives
\[
 2^{\beta_*}=2^{\alpha_2}3^{\alpha_3}W^r.
\]
The defining equation for $\rho_*$ is exactly
\[
 (r\theta-s)\rho_*
 =\gamma_2+(\gamma_3-\alpha_2)\theta-\alpha_3\theta^2,
\]
which rearranges to
\[
 \beta_*\theta=\gamma_2+\gamma_3\theta+s\rho_*.
\]
Exponentiating base $2$ gives the second equality.
\end{proof}

\begin{remark}
Theorem~\ref{thm:rank3-special-radical} does not solve the branch: deciding
whether \eqref{eq:rank3-special-radical} can be an integer is another form of
the same transcendence obstruction.  Its value is precision.  It removes
$M,N,W,$ and $\beta$ as independent unknowns and isolates one explicitly
parameterized radical whose integrality is necessary and sufficient.
In the no-cross-prime families it collapses to
$W^{|r\theta-s|}=3^d$.
\end{remark}

\section{Closure attempts and what they actually require}
\label{sec:closure-attempts}

The new invariants sharpen both branches, but neither is a concealed proof of
the conjecture.  This section records the strongest consequences that can be
extracted without importing an unproved transcendence statement, and it
states branch-specific finishing theorems precisely enough to guide the next
round of work.

\subsection{Why the energy gap does not close rank four}

Let $D_T$ denote a nonzero integer determinant in the mixed-interpolation
construction of H3.  Write its full guaranteed common divisor as
\begin{equation}
 Q_T^{\act}=Q_T^{\trop}C_T,
 \qquad C_T\in\Z_{\ge1},
 \label{eq:actual-divisor-cancellation}
\end{equation}
where $C_T$ is the gain from cancellation among terms that have the same
minimal valuation.  The coarse archimedean estimate in the inherited method
has the form
\begin{equation}
 \limsup_{T\to\infty}\frac{\log|D_T|}{H_T}\le1.
 \label{eq:coarse-archimedean-one}
\end{equation}
A contradiction would follow if the divisor exponent exceeded the right-hand
side.  By Theorem~\ref{thm:exact-cross-coprime-ceiling}, a sufficient
cancellation theorem is therefore
\begin{equation}
 \liminf_{T\to\infty}\frac{\log C_T}{H_T}
 >\frac3{10}+\frac{\E(M,N)}{20}.
 \label{eq:exact-cancellation-tax}
\end{equation}

The energy correction thus raises the cancellation tax.  This is not a defect
in the calculation: a factorization that is more visibly rank four creates
more primewise disagreement and lowers the termwise minimum.  Any successful
rank-four determinant must exploit the same disagreement a second time,
through signed cancellation, an exterior-algebraic alternation, or a sharper
archimedean norm.

More generally, if a new determinant family has archimedean exponent
$A_\infty$ and tropical divisor exponent
$7/10-\E/20$, it suffices to prove
\begin{equation}
 \liminf\frac{\log C_T}{H_T}
 > A_\infty-\frac7{10}+\frac\E{20}.
 \label{eq:general-cancellation-tax}
\end{equation}
This formula is a useful accounting identity: every proposed determinant can
be evaluated by the three numbers $A_\infty$, $\E$, and the cancellation gain.

\subsection{Why the integer minor gives no uniform analytic gap}

The minor bound is discrete, but normalization divides it by $\beta^2$:
\[
 u_pv_q-u_qv_p
 =\frac{\log p\log q}{\beta^2\log2\log3}\omega_{pq}.
\]
Consequently the forced energy separation is of order $\beta^{-4}$.  Since a
hypothetical least exponent is not known to have an absolute upper bound,
this does not produce a uniform improvement over $7/10$.  In particular,
compactness arguments on the probability simplex cannot replace the integer
minor: the arithmetic points may approach the diagonal as the valuation
heights grow.

This also explains why replacing the Gini norm by a smooth strictly convex
norm is unlikely to be enough.  Strict convexity detects noncollinearity, but
without a height-sensitive arithmetic lower bound it supplies only an
unspecified positive defect.  The normalized minor provides the correct
height scale and shows exactly how small that defect is allowed to be.

\subsection{Why $|\Delta|\ge1$ does not close rank three}

The rank-three certificate gives
\[
 1\le|\Delta|\ll\beta^3.
\]
There is no conflict between these inequalities.  To turn integrality into a
contradiction one would need an independent estimate
$|\Delta|<1$, or a nonzero algebraic quantity divisible by $\Delta$ whose
archimedean size tends to zero.  The present conic identity supplies neither:
$\Delta$ is a coefficient determinant, not the value of a small linear form.

The same issue appears in M\"obius language.  Equation
\eqref{eq:mobius-beta} places $\beta$ in the integral
$\operatorname{PGL}_2$-orbit of $\theta$, with determinant $-\Delta$.
Integral M\"obius transformations of arbitrarily large determinant are
plentiful.  The missing theorem must use the simultaneous integrality of
$2^\beta$ and $3^\beta$, not merely the orbit relation.

\subsection{Classical linear forms and exponentials theorems}

Baker's theorem is decisive for a vanishing \emph{linear} form in logarithms
of algebraic numbers with algebraic coefficients.  The rank-three relation
\eqref{eq:formal-q} is instead quadratic in $\theta$, or equivalently bilinear
in the three logarithms $X,Y,Z$.  When $\Delta=0$ the quadratic factors and
Baker-type rank reduction applies; Proposition~\ref{thm:Delta-nonzero} shows
that a counterexample is forced into the complementary nonfactorable case.
This is a structural explanation for the repeated Baker dead end in the
unified report.

The six exponentials theorem also stops one step short.  In the localized
square \eqref{eq:four-localized}, the base logarithms $\log2$ and $\log W$
span only a two-dimensional $\Q$-space.  The theorem needs three independent
base logarithms against two independent exponents.  Adjoining a third
algebraic base guarantees that at least one of six values is transcendental,
but does not force the four already-known entries in the $2\times2$ square to
contain a transcendental number.

The matrix coefficient conjecture is designed precisely for matrices of
logarithms at this rank boundary.  The current Dasgupta--Kakde framework gives
substantial structural results but does not, in the form presently available,
prove the restricted coefficient statement needed here
\cite{DasguptaKakde2026}.  Thus invoking that framework without an additional
special-case argument would only rename the remaining gap.

\subsection{Branch-specific finishing theorems}

The dichotomy suggests two sharply different targets.

\begin{conjecture}[Rank-four wedge-cancellation target]
\label{conj:wedge-cancellation-target}
For the mixed-interpolation determinant attached to a cross-coprime output
pair, the cancellation factor in \eqref{eq:actual-divisor-cancellation}
satisfies
\begin{equation}
 \liminf_{T\to\infty}\frac{\log C_T}{H_T}
 >\frac3{10}+\frac{\E(M,N)}{20}.
 \label{eq:wedge-cancellation-target}
\end{equation}
More ambitiously, the excess should admit a positive lower bound depending on
one normalized nonzero exterior minor $w_{pq}$.
\end{conjecture}

A proof of Conjecture~\ref{conj:wedge-cancellation-target}, together with
nonvanishing of the determinant, would eliminate the cross-coprime rank-four
branch.  The theorem should not be sought in termwise valuations; those have
already been optimized by H3.  It must analyze the initial form in the graded
$p$-adic algebra and show that the coefficients indexed by a nonzero wedge do
not all cancel at the tropical level.

\begin{conjecture}[Rank-three special-radical target]
\label{conj:rank3-radical-target}
For every nonnegative integral sextuple
$(\alpha_2,\alpha_3,\gamma_2,\gamma_3,r,s)$ satisfying the inherited
leastness conditions, $(r,s)\ne(0,0)$, $r\theta\ne s$, and $\Delta\ne0$, the
number in \eqref{eq:rank3-special-radical} is not an integer $>1$ coprime to
$6$ whenever the associated $\rho_*,\beta_*$ are positive.
\end{conjecture}

By Theorem~\ref{thm:rank3-special-radical}, this target eliminates the whole
rank-three branch, not merely the no-cross-prime subbranch.  Its simplest
visible special case is the following.

\begin{conjecture}[Localized no-cross-prime radical target]
\label{conj:localized-radical-target}
Let $r,s,d$ be positive integers and let $W>1$ be an integer coprime to $6$.
Then
\begin{equation}
 W^{|r\theta-s|}\ne3^d.
 \label{eq:localized-radical-target}
\end{equation}
\end{conjecture}

Conjecture~\ref{conj:localized-radical-target} is equivalent to excluding the
two families of Proposition~\ref{thm:no-cross-classification}.  It is stronger
than merely asserting that $\theta$ is irrational, but substantially more
specialized than the full four exponentials conjecture: one diagonal of the
associated exponential square is fixed to the rational integer $3$, while the
second base is constrained to be an integer coprime to $6$.

\begin{resultbox}
\textbf{Status after the new analysis.}
The output space is now partitioned by two mutually exclusive positive
integer certificates.  Rank four has an exact probability-energy loss and a
quantitative $\beta^{-4}$ separation.  Rank three has an exact integral
M\"obius discriminant, a nondegenerate formal conic, and a necessary-and-
sufficient one-radical formulation.  A complete proof still requires one new
input of the kind stated in
Conjecture~\ref{conj:wedge-cancellation-target} or
Conjecture~\ref{conj:rank3-radical-target}; neither follows from the classical
theorems presently in the ledger.
\end{resultbox}

\section{A Lean formalization blueprint}
\label{app:formalization}

The new arguments are unusually well suited to formalization.  Their hard
inputs are elementary: finite sums, unique factorization, inequalities for
nonnegative reals, $2\times2$ determinants, and polynomial identities.  The
transcendence barrier enters only at the final explicitly marked targets.
This section gives a dependency-aware Lean 4 plan.  The declarations below
are specifications rather than a claim that the code has already been
accepted by Lean.

\subsection{Representation choices}

The external valuation vectors should be represented by finitely supported
functions
\[
 a,b:\{p:\N\mid p\text{ prime},\ p\ge5\}\to_0\N.
\]
In Lean this can be implemented either by \code{Finsupp} on the subtype of external
primes or, during the analytic development, by an arbitrary finite index type.
The latter makes the probability and energy lemmas reusable and isolates all
number-theoretic bookkeeping in a thin adapter layer.

The recommended split is:
\begin{enumerate}[label=\textbf{L\arabic*.}]
\item prove the energy inequalities for a finite type $\iota$ and two
  probability vectors $u,v:\iota\to\R$;
\item prove the straddling lemma for a finite weighted family of nonzero real
  numbers with weighted sum zero;
\item instantiate $\iota$ by the finite union of the prime supports of $M$ and
  $N$;
\item prove the valuation-minor identities in \code{\Z}, coercing to \code{\R} only at
  the final normalized-wedge step;
\item formalize rank three in a separate record with six natural-number
  parameters and one positive real $\theta$.
\end{enumerate}

This organization avoids mixing \code{Nat.factorization}, real logarithms, and
finite-probability algebra in every proof.

\subsection{Proposed file layout}

\begin{table}[htbp]
\centering
\caption{Suggested Lean module decomposition.}
\label{tab:lean-modules}
\small
\begin{tabularx}{\textwidth}{@{}p{0.35\textwidth}X@{}}
\toprule
\textbf{module} & \textbf{responsibility} \\
\midrule
\code{AlaogluErdos/Theta.lean}
 & define $\theta=\log_2 3$; prove positivity, irrationality, and
   $\theta\log2=\log3$ \\
\code{AlaogluErdos/}\newline\code{FiniteProbability.lean}
 & probability vectors, total variation, diagonal energy, and
   Lemma~\ref{lem:energy-tv} \\
\code{AlaogluErdos/GiniQuadrant.lean}
 & closed first-quadrant Gini formula and the exact energy decomposition \\
\code{AlaogluErdos/PrimeSlope.lean}
 & weighted-zero-sum straddling, local nonvanishing, and positivity of an
   integral minor \\
\code{AlaogluErdos/RankFourGap.lean}
 & normalized wedge, the constant $\kappa$, and
   Theorem~\ref{thm:explicit-energy-gap} \\
\code{AlaogluErdos/ExternalRay.lean}
 & primitive-ray decomposition of proportional nonnegative \code{Finsupp} vectors \\
\code{AlaogluErdos/MobiusData.lean}
 & $U,V,\Delta$, M\"obius formula, determinant identity, and zero patterns \\
\code{AlaogluErdos/FormalConic.lean}
 & polynomial $\Phi$, symmetric matrix $S$, and $\det S=-\Delta/4$ \\
\code{AlaogluErdos/}\newline\code{SpecialRadical.lean}
 & Theorem~\ref{thm:rank3-special-radical} and no-cross-prime specialization \\
\code{AlaogluErdos/}\newline\code{BranchDichotomy.lean}
 & exterior bivector and Theorem~\ref{thm:exterior-mobius-dichotomy} \\
\bottomrule
\end{tabularx}
\end{table}

\subsection{Abstract probability layer}

A practical definition avoids division by zero explicitly.

\begin{lstlisting}[language={},caption={Lean-style specification of the
energy and total variation.}]
def diagTerm (x y : Real) : Real :=
  if max x y = 0 then 0 else (x - y)^2 / max x y

def diagEnergy [Fintype ι] (u v : ι → Real) : Real :=
  ∑ i, diagTerm (u i) (v i)

def totalVariation [Fintype ι] (u v : ι → Real) : Real :=
  (1 / 2 : Real) * ∑ i, |u i - v i|

structure ProbVec (ι : Type*) [Fintype ι] where
  val    : ι → Real
  nonneg : ∀ i, 0 ≤ val i
  sum_eq : ∑ i, val i = 1
\end{lstlisting}

The central abstract theorem should be stated independently of primes.

\begin{lstlisting}[language={},caption={Target statement for the
energy--variation comparison.}]
theorem energy_tv_bounds [Fintype ι]
    (u v : ProbVec ι) :
    4 * totalVariation u.val v.val ^ 2 /
        (1 + totalVariation u.val v.val)
      ≤ diagEnergy u.val v.val
    ∧ diagEnergy u.val v.val
      ≤ 2 * totalVariation u.val v.val := by
  -- Set d i := u i - v i and m i := max (u i) (v i).
  -- Prove ∑ i, m i = 1 + TV.
  -- Apply finite Cauchy--Schwarz to |d i| / sqrt (m i).
  -- Use |d i| ≤ m i for the upper bound.
  ...
\end{lstlisting}

A robust formal proof can avoid square roots and divisions by invoking the
Engel form of Cauchy--Schwarz on the support where $m_i>0$.  Alternatively,
prove
\[
 \left(\sum |d_i|\right)^2
 \le\left(\sum d_i^2/m_i\right)\left(\sum m_i\right)
\]
by \code{Finset.sum_mul_sq_le_sq_mul_sq} after defining zero-safe terms.

The first-quadrant Gini formula should be formalized algebraically rather than
through measure integration, because H3 already supplies the closed formula.

\begin{lstlisting}[language={},caption={Algebraic first-quadrant identity.}]
theorem gini_firstQuadrant (u v : Real) (hu : 0 ≤ u) (hv : 0 ≤ v) :
    gini u v = (2 / 15 : Real) *
      (u + v + (u - v)^2 / max u v) := by
  rcases le_total u v with huv | hvu
  · rw [gini_closed_of_nonneg_of_le hu huv]
    field_simp [max_eq_right huv]
    ring
  · rw [gini_symm, gini_closed_of_nonneg_of_le hv hvu]
    field_simp [max_eq_left hvu]
    ring
\end{lstlisting}

A practical implementation should split off the origin before
\code{field_simp}; although Lean defines division by zero to be zero, the
nonzero branches are clearer when their denominator hypotheses are explicit.

\subsection{Straddling and the integer minor}

The sign argument has no number theory in it.

\begin{lstlisting}[language={},caption={Abstract weighted straddling lemma.}]
theorem exists_pos_and_neg_of_weighted_sum_zero
    [Fintype ι] [Nonempty ι]
    (w δ : ι → Real)
    (hw : ∀ i, 0 < w i)
    (hδ : ∀ i, δ i ≠ 0)
    (hsum : ∑ i, w i * δ i = 0) :
    (∃ p, 0 < δ p) ∧ (∃ q, δ q < 0) := by
  constructor
  · by_contra h
    push_neg at h
    have : ∑ i, w i * δ i < 0 := by
      -- every δ i is strictly negative
      ...
    linarith
  · by_contra h
    push_neg at h
    have : 0 < ∑ i, w i * δ i := by
      -- every δ i is strictly positive
      ...
    linarith
\end{lstlisting}

The local nonvanishing lemma can be based on irrationality of $\theta$.

\begin{lstlisting}[language={},caption={Local slope nonvanishing and minor
positivity.}]
theorem nat_mul_theta_ne_nat
    (hθ : Irrational θ) (a b : Nat) (hab : a ≠ 0 ∨ b ≠ 0) :
    (a : Real) * θ - b ≠ 0 := by
  intro h
  have ha : a ≠ 0 := by
    intro ha0
    simp [ha0] at h
    exact hab.resolve_right (Nat.cast_eq_zero.mp h)
  have : θ = (b : Real) / a := by
    field_simp [ha]
    linarith
  exact hθ.ne_rat this

theorem minor_pos_of_straddles
    (aP bP aQ bQ : Nat)
    (hp : 0 < (aP : Real) * θ - bP)
    (hq : (aQ : Real) * θ - bQ < 0) :
    0 < (aP * bQ : Int) - aQ * bP := by
  -- Rewrite the real coercion as
  -- aP*(bQ-aQ*θ) + aQ*(aP*θ-bP), then use integrality.
  ...
\end{lstlisting}

After this theorem, the lower bound $\omega_{pq}\ge1$ is immediate from
\code{Int.add_one_le_iff} or \code{Int.one_le_iff_ne_zero} plus positivity.
The normalized wedge identity is a single \code{ring} proof once the definitions of
$u_p$ and $v_p$ are unfolded.

The inequality $w_{pq}\le\TV$ is again abstract and should be proved for any
pair of probability vectors.  This keeps all real-analysis details out of the
valuation module.

\subsection{Rank-three data as an integral record}

Subtraction makes \code{Nat} cumbersome, so the derived quantities $U,V,\Delta$
should live in \code{Int} from the outset.

\begin{lstlisting}[language={},caption={Proposed rank-three parameter record.}]
structure RankThreeData where
  a2 a3 g2 g3 r s : Nat
  rs_ne_zero : r ≠ 0 ∨ s ≠ 0

def RankThreeData.U (d : RankThreeData) : Int :=
  d.s * d.a2 - d.r * d.g2

def RankThreeData.V (d : RankThreeData) : Int :=
  d.s * d.a3 - d.r * d.g3

def RankThreeData.Delta (d : RankThreeData) : Int :=
  -(d.r : Int) * d.U - (d.s : Int) * d.V
\end{lstlisting}

The first key identity is purely integral.

\begin{lstlisting}[language={},caption={Expanded discriminant identity.}]
theorem delta_expanded (d : RankThreeData) :
    d.Delta =
      (d.r : Int)^2 * d.g2
      + (d.r : Int) * d.s * (d.g3 - d.a2)
      - (d.s : Int)^2 * d.a3 := by
  simp [RankThreeData.Delta, RankThreeData.U, RankThreeData.V]
  ring
\end{lstlisting}

The M\"obius determinant can use \code{Matrix (Fin 2) (Fin 2) Int}; for a
$2\times2$ matrix, it may be simpler to prove the explicit determinant formula
than to unfold the generic permutation definition repeatedly.

\begin{lstlisting}[language={},caption={M\"obius determinant identity.}]
def mobiusMatrix (d : RankThreeData) : Matrix (Fin 2) (Fin 2) Int :=
  !![d.V, d.U; -(d.r : Int), d.s]

theorem det_mobiusMatrix (d : RankThreeData) :
    Matrix.det (mobiusMatrix d) = -d.Delta := by
  simp [mobiusMatrix, RankThreeData.Delta]
  ring
\end{lstlisting}

The implication $\Delta=0\Rightarrow\beta\in\Q$ should be split by the zero
patterns $r=0$, $s=0$, and $r,s>0$, exactly as in
Proposition~\ref{thm:Delta-nonzero}.  This is longer than the paper proof but
avoids hidden denominator assumptions.  The final contradiction can use the
already formalized rational-exponent lemma:
\[
 x\in\Q,\quad 2^x\in\Z \quad\Longrightarrow\quad x\in\Z.
\]

\subsection{Formal quadratic bridge}

The identity $\det S=-\Delta/4$ is ideal for \code{ring_nf}.  To avoid fractions,
formalize the doubled matrix
\begin{equation}
 2S=
 \begin{pmatrix}
 2\gamma_2&\gamma_3-\alpha_2&s\\
 \gamma_3-\alpha_2&-2\alpha_3&-r\\
 s&-r&0
 \end{pmatrix},
 \label{eq:doubled-S}
\end{equation}
for which
\begin{equation}
 \det(2S)=-2\Delta.
 \label{eq:doubled-S-det}
\end{equation}
This stays entirely in \code{Int}; the rational identity follows by coercion.

\begin{lstlisting}[language={},caption={Integer version of the formal-conic
determinant.}]
def doubledConicMatrix (d : RankThreeData) : Matrix (Fin 3) (Fin 3) Int :=
  !![2*d.g2, d.g3-d.a2, d.s;
     d.g3-d.a2, -2*d.a3, -d.r;
     d.s, -d.r, 0]

theorem det_doubledConicMatrix (d : RankThreeData) :
    Matrix.det (doubledConicMatrix d) = -2 * d.Delta := by
  simp [doubledConicMatrix, RankThreeData.Delta,
        RankThreeData.U, RankThreeData.V, Matrix.det_fin_three]
  ring
\end{lstlisting}

The conic identity itself is best stated in a commutative ring, not only in
$\R$.  That makes explicit that no analytic property of logarithms is used in
the bridge.

\subsection{Special-radical equivalence at the logarithmic level}

Formalizing \code{Real.rpow} too early creates avoidable side conditions.  First
prove an equivalence between two affine equations in $\beta,\rho$ and the
single formula for $\rho_*$.  Only then instantiate
$\rho=\log W/\log2$ and exponentiate.

\begin{lstlisting}[language={},caption={Linear-algebra core of the special
radical reduction.}]
theorem rho_eq_iff_two_log_equations
    (d : RankThreeData) (hden : (d.r : Real) * θ ≠ d.s)
    (rho beta : Real) :
    (beta = d.a2 + d.a3 * θ + d.r * rho ∧
     beta * θ = d.g2 + d.g3 * θ + d.s * rho) ↔
    (rho =
      (d.g2 + (d.g3-d.a2)*θ - d.a3*θ^2) /
      ((d.r : Real)*θ-d.s) ∧
     beta = d.a2 + d.a3*θ + d.r*rho) := by
  constructor <;> intro h
  · rcases h with ⟨h1,h2⟩
    constructor
    · field_simp [hden]
      nlinarith
    · exact h1
  · rcases h with ⟨hr,hb⟩
    constructor
    · exact hb
    · field_simp [hden] at hr
      nlinarith
\end{lstlisting}

The final exponential statements use the standard identities
\[
 \code{Real.rpow_add},\qquad
 \code{Real.rpow_natCast},\qquad
 2^\theta=3.
\]
Because all bases are positive, no branch-cut issue occurs.

\subsection{Dichotomy and dependency graph}

The exterior dichotomy is a finite-rank statement about \code{Finsupp} vectors.
Mathlib already provides exterior algebra, but importing it is unnecessary for
this theorem: define \code{allMinorsZero a b} and prove directly that all minors
vanish iff the two nonnegative integral vectors lie on a common primitive ray.
The actual bivector can be added later as a conceptual wrapper.

\begin{center}
\begin{tabularx}{0.96\textwidth}{@{}>{\centering\arraybackslash}X>{\centering\arraybackslash}X@{}}
\toprule
\textbf{rank-four dependency chain} & \textbf{rank-three dependency chain} \\
\midrule
\code{Theta} $\longrightarrow$ \code{PrimeSlope}
& \code{Theta} $\longrightarrow$ \code{MobiusData} \\
\code{FiniteProb.} $\longrightarrow$ \code{RankFourGap}
& \code{ExternalRay} $\longrightarrow$ \code{MobiusData} \\
\code{PrimeSlope} $\longrightarrow$ \code{RankFourGap}
& \code{MobiusData} $\longrightarrow$ \code{FormalConic} \\
\code{RankFourGap} $\searrow$
& \code{MobiusData} $\longrightarrow$ \code{SpecialRadical} \\
\multicolumn{2}{c}{\rule{0pt}{2.4ex}$\phantom{X}\longrightarrow$ \code{BranchDichotomy} $\longleftarrow\phantom{X}$} \\
\bottomrule
\end{tabularx}
\end{center}

A useful formalization milestone is the following theorem with all analytic
hypotheses abstracted away:

\begin{quote}
Given two nonnegative integral external valuation vectors of positive total
logarithmic masses satisfying the slope-conservation identity, either a
positive integral minor exists and forces the normalized energy bound, or the
vectors have a primitive common ray and the associated M\"obius determinant
is nonzero whenever the induced exponent is irrational.
\end{quote}

This theorem captures every new unconditional argument in the report.  The two
finishing conjectures should remain separate axioms or theorem holes; hiding
either inside a broad transcendence import would defeat the diagnostic value
of the formalization.

\section{Symbolic, numerical, and logical audit}
\label{sec:audit}

The formulas in this report were checked in three complementary ways:
hand derivation, exact symbolic expansion, and limiting-case tests.  This
section records enough detail for independent reproduction.

\subsection{Exact polynomial identities}

The following compact SymPy session verifies the three identities most prone
to sign or factor errors.

\begin{lstlisting}[language=Python,caption={Exact symbolic checks used in the
report.}]
import sympy as s

u, v = s.symbols('u v', positive=True)
g0 = 2*(2*u**2 + 3*u*(v-u) + 2*(v-u)**2)/(15*v)
g1 = s.Rational(2,15)*(u+v+(u-v)**2/v)
assert s.simplify(g0-g1) == 0          # first-quadrant Gini

a2,a3,g2,g3,r,t = s.symbols('a2 a3 g2 g3 r s')
Delta = r**2*g2 + r*t*(g3-a2) - t**2*a3
S = s.Matrix([
 [g2,       (g3-a2)/2,  t/2],
 [(g3-a2)/2,   -a3,    -r/2],
 [t/2,          -r/2,     0]])
assert s.factor(S.det()) == -Delta/4   # formal conic

U = t*a2-r*g2
V = t*a3-r*g3
B = s.Matrix([[V,U],[-r,t]])
assert s.factor(B.det()) == -Delta     # Mobius determinant
\end{lstlisting}

The doubled integral matrix \eqref{eq:doubled-S} provides an independent
factor check: because a $3\times3$ determinant is cubic under scalar
multiplication,
\[
 \det(2S)=2^3\det S=8\left(-\frac\Delta4\right)=-2\Delta.
\]

\subsection{Numerical constants}

Using natural logarithms,
\begin{align}
 \kappa
 &=\frac{\log5\log7}{\log2\log3}
   =4.1127006240\ldots,\\
 \frac{\kappa^2}{5}
 &=3.3828612847\ldots,\\
 \frac{\kappa^2}{10}
 &=1.6914306423\ldots .
 \label{eq:numerical-kappa-data}
\end{align}
The coefficient in \eqref{eq:simple-quantitative-ceiling} is the last of these
numbers.  The constants are dimensionless and independent of the logarithm
base.

The cubic discriminant bound uses
\begin{equation}
 \frac{3\log2\log3}{(\log5)^2}
 =0.8819474591\ldots .
 \label{eq:numerical-delta-constant}
\end{equation}

\subsection{Limiting cases for the energy}

Three test configurations expose the normalization and all endpoint factors.
They are tests of the analytic identity, not claims that the configurations
come from counterexamples.

\paragraph{Coincident distributions.}
If $u=v$, then $\E=0$ and
\eqref{eq:exact-cross-coprime-ceiling} gives $7/10$.  Theorem
\ref{thm:exact-cross-coprime-ceiling} shows that arithmetic factorizations
cannot realize this endpoint.

\paragraph{Disjoint distributions.}
For $u=(1,0)$ and $v=(0,1)$, one has $\TV=1$ and $\E=2$; hence the tropical
constant is $7/10-2/20=3/5$.  Both inequalities in the allowed range are
therefore sharp on the abstract probability simplex.

\paragraph{Symmetric near-diagonal perturbation.}
For $0<\varepsilon<1/2$, let
\[
 u=\left(\frac12+\varepsilon,\frac12-\varepsilon\right),
 \qquad
 v=\left(\frac12-\varepsilon,\frac12+\varepsilon\right).
\]
Then
\begin{equation}
 \TV=2\varepsilon,
 \qquad
 \E=\frac{16\varepsilon^2}{1+2\varepsilon}.
 \label{eq:near-diagonal-test}
\end{equation}
Thus the lower inequality in Lemma~\ref{lem:energy-tv} is attained exactly.
This verifies both the factor $4$ and the denominator $1+\TV$ and shows that
the quadratic small-defect scale is optimal.

\subsection{Sign tests for the M\"obius discriminant}

The axis cases provide immediate sign checks:
\[
 r=0<s \Longrightarrow \Delta=-s^2\alpha_3,
 \qquad
 s=0<r \Longrightarrow \Delta=r^2\gamma_2.
\]
These agree with both $-rU-sV$ and the expanded quadratic formula.

In no-cross family $(+)$, substituting
$\alpha_2=\alpha_3=\gamma_2=0$, $\gamma_3=d$ gives
$\Delta=rsd>0$.  In family $(-)$, substituting
$\alpha_3=\gamma_2=\gamma_3=0$, $\alpha_2=d$ gives
$\Delta=-rsd<0$.  The signs agree with the direction of the approximation
$r\theta-s$.

A useful synthetic unit-determinant pattern is
\[
 \alpha_2=\alpha_3=\gamma_3=0,
 \quad \gamma_2=1,
 \quad r=s=1,
\]
for which $\Delta=1$ and
\[
 \rho_*=\beta_*=\frac1{\theta-1},
 \qquad
 W_*=2^{1/(\theta-1)}.
\]
Nothing elementary forces this special radical to be nonintegral.  This test
confirms that $|\Delta|=1$ is a genuine boundary case rather than an immediate
contradiction.

\subsection{Dependency and noncircularity audit}

The unconditional chain in rank four is
\[
 \begin{aligned}
 \text{factorizations}
 &\Longrightarrow \text{probability vectors}
 \Longrightarrow \text{slope conservation}
 \Longrightarrow \omega_{pq}\ge1,\\
 &\Longrightarrow \TV\ge\kappa/\beta^2
 \Longrightarrow \E\ge
 \frac{4\kappa^2}{\beta^4+\kappa\beta^2}.
 \end{aligned}
\]
No determinant nonvanishing conjecture or four-exponentials statement is used
in this chain.  H3 is used only in the final conversion from energy to the
tropical divisor exponent.

The unconditional chain in rank three is
\[
 A\wedge B=0
 \Longrightarrow (W,r,s)
 \Longrightarrow \beta=\frac{U+V\theta}{s-r\theta}
 \Longrightarrow \Delta\ne0
 \Longrightarrow \det S=-\Delta/4
 \Longrightarrow \text{special radical }W_*.
\]
The proof that $\Delta\ne0$ uses only the inherited rational-exponent lemma and
the assumption that $\beta$ is a counterexample.  The localized
four-exponentials square is a consequence, not an input.

The public model-challenge reports and current literature reconnaissance are
contextual only.  No mathematical step in the two chains depends on a claimed
external breakthrough, unpublished formalization, or machine-generated proof.

\section{Conclusions and research priorities}
\label{sec:conclusion}

This continuation does not claim a complete proof of the Alaoglu--Erdős
conjecture.  It does produce a branch-exhaustive arithmetic refinement that
was absent from the unified report.

In rank four, the normalized external factorizations are probability
measures.  The first-quadrant Gini norm splits exactly into total mass plus a
max-$\chi^2$ energy.  Unique factorization excludes zero energy, and a signed
prime-slope conservation law forces a positive integer valuation minor.
That minor yields an explicit total-variation separation and the quantitative
ceiling
\[
 \frac7{10}-
 \frac{\kappa^2}{5(\beta^4+\kappa\beta^2)}.
\]
The old strictness clause is therefore replaced by an exact formula and a
height-sensitive arithmetic deficit.

In rank three, vanishing of every external minor produces a primitive core
$W$.  The remaining structural information is measured by
\[
 \Delta=-rU-sV=\det(-B),
\]
a nonzero integer.  The same number is, up to $-1/4$, the determinant of the
formal quadratic relation among $\log2,\log3,$ and $\log W$.  It classifies the
zero patterns, separates the cross-coprime branch, yields a cubic height
bound, and turns the whole rank-three branch into the integrality question for
one explicit special radical.

The most promising next work is correspondingly bifurcated.

\begin{enumerate}[label=\textbf{P\arabic*.},leftmargin=3em]
\item \textbf{Compute initial forms of the mixed determinant.}
  The relevant object is not the minimum valuation but the sum of determinant
  terms attaining that minimum.  The integer wedge $\omega_{pq}$ should be
  inserted into this graded calculation from the start.  The concrete target
  is \eqref{eq:wedge-cancellation-target}.

\item \textbf{Search for a rank-three special-radical theorem.}
  Work directly with \eqref{eq:rank3-special-radical}, exploiting
  nonnegativity, the leastness condition
  $\min(\alpha_2,\gamma_3)=0$, coprimality of $W$ with $6$, and
  $|\Delta|\ge1$.  These restrictions are invisible in a generic
  four-exponentials formulation and may permit a specialized $S$-unit,
  modular, or $p$-adic argument.

\item \textbf{Formalize the unconditional core.}
  The energy inequality, straddling minor, M\"obius determinant, and formal
  conic are small, independent Lean milestones.  A formal development would
  make sign conventions and zero cases completely stable while leaving the
  two genuine transcendence holes exposed.

\item \textbf{Use computation only as a theorem-discovery tool.}
  Enumerate bounded rank-three parameter sextuples, compute $\Delta$ and
  high-precision $\rho_*$, and test proximity of $2^{\rho_*}$ to integers.
  Such a search cannot prove nonintegrality, but exceptional families may
  reveal congruence obstructions or a missing descent invariant.
\end{enumerate}

The main conceptual gain is that the two hard branches now fail for different,
explicit reasons.  Rank four lacks enough cancellation in an otherwise
quantified determinant.  Rank three lacks an integrality obstruction for an
otherwise explicit M\"obius special radical.  Treating those as separate
problems is more informative than another global reformulation of the original
$2\times2$ exponential square.

\appendix

\section{Compact theorem ledger}
\label{app:theorem-ledger}

\begin{longtable}{@{}p{0.24\textwidth}p{0.18\textwidth}p{0.50\textwidth}@{}}
\caption{New statements and their logical status.}\label{tab:new-ledger}\\
\toprule
\textbf{statement} & \textbf{status} & \textbf{content} \\
\midrule
\endfirsthead
\toprule
\textbf{statement} & \textbf{status} & \textbf{content} \\
\midrule
\endhead
Proposition~\ref{prop:gini-first-quadrant}
 & unconditional
 & exact first-quadrant Gini formula with diagonal-defect term \\
Theorem~\ref{thm:exact-cross-coprime-ceiling}
 & unconditional from H3
 & tropical exponent $7/10-\E/20$ and exclusion of equality \\
Lemma~\ref{lem:prime-slope-straddling}
 & unconditional
 & opposite local slopes force a positive integral valuation minor \\
Theorem~\ref{thm:explicit-energy-gap}
 & unconditional from H1--H3
 & $\TV\ge\kappa/\beta^2$ and explicit $\beta^{-4}$ energy gap \\
Proposition~\ref{thm:Delta-nonzero}
 & unconditional from H1--H4
 & rank-three M\"obius determinant is a nonzero integer \\
Proposition~\ref{thm:determinant-bridge}
 & unconditional
 & conic matrix determinant is $-\Delta/4$ \\
Proposition~\ref{thm:no-cross-classification}
 & unconditional from H1--H4
 & two rank-three no-cross-prime families and
   $W^{|r\theta-s|}=3^d$ \\
Theorem~\ref{thm:localized-square}
 & unconditional
 & second explicit four-exponentials square in the rank-three no-cross branch \\
Theorem~\ref{thm:exterior-mobius-dichotomy}
 & unconditional from H4
 & exhaustive split by $\Omega$; active integer certificate $I_4$ or $I_3$ \\
Theorem~\ref{thm:rank3-special-radical}
 & unconditional algebra
 & necessary-and-sufficient one-radical reduction of rank three \\
Conjecture~\ref{conj:wedge-cancellation-target}
 & open target
 & sufficient cancellation theorem for the cross-coprime rank-four branch \\
Conjecture~\ref{conj:rank3-radical-target}
 & open target
 & sufficient special-radical exclusion for all of rank three \\
\bottomrule
\end{longtable}

\section{Formula sheet}
\label{app:formula-sheet}

For quick reference, the two branches can be summarized as follows.

\subsection*{Rank four}
\begin{align*}
 u_p&=\frac{a_p\log p}{\beta\log2},
 &v_p&=\frac{b_p\log p}{\beta\log3},
 &\sum u_p&=\sum v_p=1,\\
 \E&=\sum_p\frac{(u_p-v_p)^2}{\max(u_p,v_p)},
 &\TV&=\frac12\sum_p|u_p-v_p|,\\
 \frac{\log Q_T^{\trop}}{H_T}
 &\longrightarrow\frac7{10}-\frac\E{20},
 &\omega_{pq}&=a_pb_q-a_qb_p\in\Z\setminus\{0\},\\
 u_pv_q-u_qv_p
 &=\frac{\log p\log q}{\beta^2\log2\log3}\omega_{pq},
 &\kappa&=\frac{\log5\log7}{\log2\log3},\\
 \TV&\ge\frac\kappa{\beta^2},
 &\E&\ge\frac{4\kappa^2}{\beta^4+\kappa\beta^2}.
\end{align*}

\subsection*{Rank three}
\begin{align*}
 M&=2^{\alpha_2}3^{\alpha_3}W^r,
 &N&=2^{\gamma_2}3^{\gamma_3}W^s,\\
 U&=s\alpha_2-r\gamma_2,
 &V&=s\alpha_3-r\gamma_3,\\
 \beta&=\frac{U+V\theta}{s-r\theta},
 &\Delta&=-rU-sV\in\Z\setminus\{0\},\\
 \Phi(X,Y,Z)
 &=\gamma_2X^2+(\gamma_3-\alpha_2)XY-\alpha_3Y^2+sXZ-rYZ,\\
 \det S&=-\frac\Delta4,
 &1&\le|\Delta|\ll\beta^3,\\
 \rho_*&=
 \frac{\gamma_2+(\gamma_3-\alpha_2)\theta-\alpha_3\theta^2}
      {r\theta-s},
 &W_*&=2^{\rho_*}.
\end{align*}

\begin{thebibliography}{99}

\bibitem{UnifiedReport}
V.~Reshetnikov and collaborating language models,
\emph{A Unified Research Report on the Alaoglu--Erdős Exponential
Integrality Conjecture}, consolidated working report, 2026.
The attachment supplied with the present continuation.

\bibitem{AlaogluErdos1944}
L.~Alaoglu and P.~Erdős,
``On highly composite and similar numbers,''
\emph{Transactions of the American Mathematical Society}
\textbf{56} (1944), 448--469.

\bibitem{Baker1966}
A.~Baker,
``Linear forms in the logarithms of algebraic numbers. I,''
\emph{Mathematika} \textbf{13} (1966), 204--216.

\bibitem{Waldschmidt2000}
M.~Waldschmidt,
\emph{Diophantine Approximation on Linear Algebraic Groups},
Grundlehren der mathematischen Wissenschaften 326,
Springer, 2000.

\bibitem{DasguptaKakde2026}
S.~Dasgupta and M.~Kakde,
``Ranks of matrices of logarithms of algebraic numbers II:
the matrix coefficient conjecture,''
arXiv:2408.08178, version current in August 2026.

\bibitem{AnthropicOpus46}
Anthropic,
``Introducing Claude Opus 4.6,'' official product announcement,
5 February 2026,
\url{https://www.anthropic.com/news/claude-opus-4-6}.

\bibitem{ByClaudeHome}
ByClaude,
``Claude Mathematical Challenges,'' public challenge index,
accessed 22 August 2026,
\url{https://byclaude.net/}.

\bibitem{ByClaudeJacobian}
ByClaude,
``The Jacobian conjecture,'' public challenge record,
accessed 22 August 2026,
\url{https://byclaude.net/talks/3}.

\bibitem{JacobianVerifier}
L.~Alp\"oge,
``Jacobian conjecture proof verifier,'' independent verification site,
accessed 22 August 2026,
\url{https://jacobian.alpoge.org/}.

\bibitem{ByClaudeHadamard}
ByClaude,
``Constructing Hadamard matrices of all orders below 2000,''
public challenge record, accessed 22 August 2026,
\url{https://byclaude.net/talks/37}.

\bibitem{ByClaudeRank30}
ByClaude,
``An elliptic curve over $\Q$ of rank at least 30,''
public challenge record, accessed 22 August 2026,
\url{https://byclaude.net/talks/47}.

\end{thebibliography}

\end{document}
